% Options for packages loaded elsewhere
\PassOptionsToPackage{unicode}{hyperref}
\PassOptionsToPackage{hyphens}{url}
%
\documentclass[
  10pt,
]{article}
\usepackage{lmodern}
\usepackage{amssymb,amsmath}
\usepackage{ifxetex,ifluatex}
\ifnum 0\ifxetex 1\fi\ifluatex 1\fi=0 % if pdftex
  \usepackage[T1]{fontenc}
  \usepackage[utf8]{inputenc}
  \usepackage{textcomp} % provide euro and other symbols
\else % if luatex or xetex
  \usepackage{unicode-math}
  \defaultfontfeatures{Scale=MatchLowercase}
  \defaultfontfeatures[\rmfamily]{Ligatures=TeX,Scale=1}
\fi
% Use upquote if available, for straight quotes in verbatim environments
\IfFileExists{upquote.sty}{\usepackage{upquote}}{}
\IfFileExists{microtype.sty}{% use microtype if available
  \usepackage[]{microtype}
  \UseMicrotypeSet[protrusion]{basicmath} % disable protrusion for tt fonts
}{}
\makeatletter
\@ifundefined{KOMAClassName}{% if non-KOMA class
  \IfFileExists{parskip.sty}{%
    \usepackage{parskip}
  }{% else
    \setlength{\parindent}{0pt}
    \setlength{\parskip}{6pt plus 2pt minus 1pt}}
}{% if KOMA class
  \KOMAoptions{parskip=half}}
\makeatother
\usepackage{xcolor}
\IfFileExists{xurl.sty}{\usepackage{xurl}}{} % add URL line breaks if available
\IfFileExists{bookmark.sty}{\usepackage{bookmark}}{\usepackage{hyperref}}
\hypersetup{
  hidelinks,
  pdfcreator={LaTeX via pandoc}}
\urlstyle{same} % disable monospaced font for URLs
\usepackage[a4paper,margin=18mm]{geometry}
\setlength{\emergencystretch}{3em} % prevent overfull lines
\providecommand{\tightlist}{%
  \setlength{\itemsep}{0pt}\setlength{\parskip}{0pt}}
\setcounter{secnumdepth}{-\maxdimen} % remove section numbering

\usepackage{esint}
\usepackage{xeCJK}
\setmainfont{DejaVu Serif}
\setmonofont{DejaVu Sans Mono}
\setCJKmainfont{Noto Serif CJK SC}
\setCJKmonofont{Noto Sans Mono CJK SC}
\setmathfont{Asana Math}
\setmathfont{Latin Modern Math}[range=\lrcorner]
\setlength{\emergencystretch}{3em}
\allowdisplaybreaks
\hypersetup{hypertexnames=false}

\author{}
\date{}

\begin{document}

\hypertarget{endpoint-geometry-actual-annuli-and-sharp-subfamilies}{%
\section{Endpoint geometry, actual annuli, and sharp
subfamilies}\label{endpoint-geometry-actual-annuli-and-sharp-subfamilies}}

Research snapshot: 844e62f7304d9f890914a6a4370ec39291ddc26e.

Complete recent proofs. The first paper's unresolved bounded-Ricci
Euler-density case is later addressed in the second paper's Section 3.
Neither version proves an unrestricted integral-Ricci endpoint theorem.
The handle-only exact infimum is not the balanced-family coefficient.

Read 00\_READ\_FIRST.md for superseded statements. Each source below is
complete and byte-preserved. Original local links remain relative to
their source; locate other texts using Volume 6. Formula numbers restart
in each source.

\begin{center}\rule{0.5\linewidth}{0.5pt}\end{center}

\hypertarget{source-1-20260927_covering_geometry_and_sharp_subfamilies.md}{%
\subsection{Source 1:
20260927\_covering\_geometry\_and\_sharp\_subfamilies.md}\label{source-1-20260927_covering_geometry_and_sharp_subfamilies.md}}

Original path:
continuation/20260927\_covering\_geometry\_and\_sharp\_subfamilies.md

SHA-256:
fdbfa874d5611fe54ad75ad00cd0a2f05946777c2598b6d4922abf09df40ad76

BEGIN COMPLETE SOURCE

\hypertarget{local-covering-geometry-sharp-surgery-costs-and-first-order-ale-matching}{%
\section{Local covering geometry, sharp surgery costs, and first-order
ALE
matching}\label{local-covering-geometry-sharp-surgery-costs-and-first-order-ale-matching}}

27 September 2026. Continuation from local commit \texttt{46e7efc}.

This manuscript continues the actual-field theory; it does not replace
the frozen research library. The inherited results used here are the
bounded actual-form supplier, the local incompatibility theorem, and the
conformal minimizing-sequence reduction in
\href{20260927_conformal_separation_and_zero_end.md}{the preceding
continuation}. The unrestricted integral zero-end problem and the best
coefficient on the balanced blowup/handle family are not assumed solved.

The first connection strengthens base collapse to collapse of
\textbf{local universal covers}, then uses verified smoothing inputs to
exclude further approximate geometries. The second determines an exact
all-metric infimum on a real subfamily and removes scalar energy from
the remaining chiral filling criterion. The third excludes a precisely
stated global ALE matching mechanism, retaining the first-order source
and then weakening pointwise decay to finite energy. Classical geometric
and gauge inputs are distinguished from these deductions.

\hypertarget{local-covering-volume-detects-the-forbidden-curvature-mode}{%
\subsection{1. Local covering volume detects the forbidden curvature
mode}\label{local-covering-volume-detects-the-forbidden-curvature-mode}}

Let a closed oriented four-manifold carry a smooth metric with
\texttt{\textbar{}Ric\textbar{}\textless{}=K}, a smooth self-dual
two-form \texttt{\textbar{}omega\textbar{}\textless{}=L}, and a constant
\texttt{t\ in\ {[}0,1/4{]}}. Set

\[
 Z_t=|\operatorname{Ric}+g|^2+
 |W^+-t(|\omega|^2 I/3-\omega\otimes\omega)|^2
 +( |\omega|^2-2)^2+|d\omega|^2.                 \tag{1.1}
\]

For a fixed radius rho\textgreater0 let

\[
 \pi_x:(\widetilde {B_{4\rho}(x)},\tilde x)
                   \longrightarrow(B_{4\rho}(x),x),\qquad
 w_\rho(x)=\rho^{-4}
       \operatorname{Vol}_{\pi_x^*g}B_\rho(\tilde x)   \tag{1.2}
\]

be the local universal cover with the lifted metric. This is not the
universal cover of the whole manifold and not an exponential chart.

\textbf{Theorem 1 (covering-volume fraction).} If G is any measurable
set of centers with \texttt{w\_rho(x)\textgreater{}=v\textgreater{}0},
then

\[
 \boxed{\frac{\operatorname{Vol}(G)}{V}
                  \le C(K,L,\rho,v)\fint_M Z_t.}      \tag{1.3}
\]

There is no bound on covering multiplicity, diameter, total volume,
sectional curvature, or the base injectivity radius.

\textbf{Proof.} The local coercivity argument of the preceding
manuscript, Theorem 6, applies on a ball of any fixed radius. It gives

\[
 \int_{B_{2\rho}(\tilde x)}\pi_x^*Z_t
                  \ge\epsilon(K,L,\rho,v)>0.         \tag{1.4}
\]

Only interior balls are used. Although the covering of an open ball need
not be complete, all balls involved here have a definite buffer from its
metric boundary; geodesics of length at most 2rho starting at the center
stay over B\_\{2rho\}(x). The interior volume comparison,
harmonic-coordinate and compactness arguments are therefore valid. One
can also repeat the proof directly on these buffered balls. On rescaling
the radius, the limiting equation is Ric=-rho\^{}2 g rather than Ric=-g;
rho is fixed, its constant is strictly negative, and the forbidden
parameter interval is rescaled by the same factor. Thus no
scale-invariance of the normalized equation is being presumed.

Let \texttt{c\_x(theta)} be the cut time on this cover. Polar
integration \textbf{before the cover cut time}, followed by Bishop's
radial Jacobian upper bound, gives

\[
 \begin{aligned}
 \int_{B_{2\rho}(\tilde x)}\pi_x^*Z_t
 &\le C(K,\rho)\int_{S_xM}
       \int_0^{\min\{2\rho,c_x(\theta)\}}
                    s^3Z_t(\exp_x(s\theta))\,ds\,d\theta\\
 &\le C(K,\rho)\int_{S_xM}\int_0^{2\rho}
                    s^3Z_t(\exp_x(s\theta))\,ds\,d\theta .
 \end{aligned}\tag{1.5}
\]

The last inequality drops the cut restriction only after replacing the
Jacobian by its upper bound; it uses nonnegativity of Z. It does not
apply Bishop comparison beyond a conjugate point. The projected geodesic
can cross the cut locus of the base without affecting the argument.

Integrate (1.4)--(1.5) over G and enlarge the nonnegative right side to
all centers. The Liouville measure on the unit tangent bundle is
preserved by the geodesic flow of the \textbf{original single metric g}.
Hence

\[
 \begin{aligned}
 \epsilon\operatorname{Vol}(G)
 &\le C(K,\rho)\int_0^{2\rho}s^3
    \int_{SM}Z_t(\exp_x(s\theta))\,d\theta\,d\mu(x)\,ds\\
 &=C(K,\rho)|S^3|\frac{(2\rho)^4}{4}\int_M Z_t.
 \end{aligned}\tag{1.6}
\]

This proves (1.3), with no counting of sheets. QED.

The needed local coercivity is itself based on a regular harmonic chart
from \href{https://arxiv.org/pdf/1406.6534}{Cheeger--Naber} and the
local incompatibility of negative Einstein ASD almost-Kaehler geometry
in
\href{https://johnarmstrongmaths.com/publications/resources/thesis.pdf}{Armstrong,
Theorem 3.3.5}. Its proof, including the t\textgreater0 Bianchi case,
remains in the preceding manuscript, Section 8. The new step here is the
cover-to-base averaging.

For a sequence with uniformly bounded Ricci and amplitude and
\texttt{avg\ Z\_t\ -\textgreater{}\ 0}, (1.3) implies, for each fixed
rho and v, that the fraction where \texttt{w\_rho\textgreater{}=v} tends
to zero in outer measure. Indeed the continuous geodesic-average right
side of (1.5) is at least epsilon at every such center, giving a
measurable majorant for that level set. Bishop comparison also bounds
w\_rho above by a constant depending only on K,rho. Splitting into
\texttt{w\_rho\textless{}v} and its complement yields convergence of the
normalized upper integral, without asserting an unproved measurability
property of the family of covers:

\[
                        \fint_M^{*} w_\rho\longrightarrow0. \tag{1.7}
\]

This is stronger than base collapse. It is not pointwise collapse at
every center and does not produce a sequence realizing such
degeneration.

\hypertarget{actual-monopoles-and-maximal-rank-collapse}{%
\subsection{2. Actual monopoles and maximal-rank
collapse}\label{actual-monopoles-and-maximal-rank-collapse}}

The inherited actual supplier, for beta\textgreater=4 and positive
projection p, gives the SAME bounded form with \texttt{L=3\ sqrt(2)},
\texttt{t=1/beta}, and

\[
 \fint Z_t\le C\left(r_{\rm Ric}+D_\beta/V\right)+o(1),
 \qquad r_{\rm Ric}=\fint|\operatorname{Ric}+g|^2.     \tag{2.1}
\]

The last allowance is freely selectable in the auxiliary construction;
the geometric set G is independent of it. Send that allowance to zero
after applying Theorem 1. No class-square condition or Hessian budget is
needed in this composition. The actual class must remain fixed through
the auxiliary weights required by the supplier.

\textbf{Corollary 2.} Under bounded Ricci,

\[
 \frac{\operatorname{Vol}(G)}V
       \le C(K,\rho,v)(r_{\rm Ric}+D_\beta/V),
                          \qquad\beta\ge4.           \tag{2.2}
\]

In particular a fixed positive volume fraction of locally noncollapsed
universal covers gives a uniform forbidden-endpoint gap, even when the
base collapses everywhere. For the original pointwise input
\texttt{delta=\textbar{}\textbar{}Ric\^{}sharp+I\textbar{}\textbar{}\_infinity,op},
\texttt{delta\textless{}=1/4}, and
\texttt{Vol(G)\textgreater{}=vartheta\ V}, this becomes

\[
             \delta^2+D_\beta/V
                   \ge\epsilon(\rho,v,\vartheta)>0.  \tag{2.3}
\]

There is a checkable topological source of these good covers.
\href{https://arxiv.org/pdf/1412.1326}{Naber--Zhang, Proposition 5.9},
printed pages 53--54, states: under Ric\textgreater=-(n-1), local
completeness, GH closeness of B\_2(p) to a specified k-dimensional
Ricci-limit ball, and

\[
 \operatorname{rank}\operatorname{Image}
       [\pi_1(B_{\delta_0}(p))\longrightarrow\pi_1(B_2(p))]=n-k,
                                                        \tag{2.4}
\]

the radius-1/2 ball in the local universal cover of B\_1(p) has a
positive volume lower bound. Their constants depend on the specified
limit ball. At n=4, after scaling a fixed radius R with
\texttt{R\^{}2\ K\textless{}=3}, this provides the cover volume needed
in (1.2) with rho=R/4. Indeed the source controls B\_\{R/2\} in the
cover of B\_R; relative volume comparison gives a lower bound for
B\_\{R/4\} in that same cover.

Consequently, a positive fraction of points satisfying these uniform
GH/model and maximal fibered-rank conditions yields (2.2)--(2.3). A
zero-end sequence cannot retain a positive fraction of any such fixed
maximal-rank local collapse geometry. If a fixed GH model controls
almost all centers, the maximal-rank fraction tends to zero; it is not
enough for the base to be collapsed. Without uniform model control one
must retain the alternative that the model constants also degenerate.

This uses the \textbf{cover-noncollapse} proposition, not the stronger
harmonic radius theorem with its more restrictive model splitting.
Neither a uniform global fundamental group nor global asphericity is
asserted to imply (2.4).

\hypertarget{a-diameter-free-upper-energy-margin-with-covering-geometry}{%
\subsection{3. A diameter-free upper-energy margin with covering
geometry}\label{a-diameter-free-upper-energy-margin-with-covering-geometry}}

Retain actual \texttt{p\textgreater{}=A\textgreater{}0}, put
\texttt{T\_+=A(1-h)},
\texttt{e=integral\ \textbar{}Ric\textbar{}\^{}2/(8\ pi\^{}2\ A)}, and

\[
 G(h)=\frac{(4+\sqrt{1-3h})^2}{9},\qquad
                        1/4\le h\le1/3.             \tag{3.1}
\]

Suppose a fixed fraction vartheta of centers has
\texttt{w\_rho\textgreater{}=v}. Then

\[
 \boxed{\delta_-^2+(G(h)-e)_+
                      \ge c(\rho,v,\vartheta)>0.}    \tag{3.2}
\]

This is a direct continuation of the precise scalar-variance calculation
in Section 13 of the preceding manuscript: replace its diameter-derived
gap by (2.3). To make the dependence transparent, decrease the constant
epsilon in (2.3) to at most 1/16. For delta\textless=1/8 and
\texttt{(G-e)\_+\textless{}=1/16}, that calculation gives

\[
 \delta^2+\inf_{\beta\ge4}D_\beta/V
           \le200\bigl(\delta+\sqrt{(G-e)_+}\bigr).
\]

Thus \texttt{c=min\{1/64,epsilon\^{}2/80000\}} suffices; the two omitted
large-error cases satisfy the same bound immediately. The calculation
keeps the total scalar variance and never replaces a near-Einstein
metric by a Yamabe metric in a pointwise argument. Since geometric
noncollapse was the sole use of diameter and class-square there, neither
a diameter bound nor \texttt{a\^{}2\textgreater{}=T\_+} nor an absolute
\texttt{p\textgreater{}=A\_0} is needed here.

For exact Ric=-g the stronger inherited optimization gives

\[
                 e\le G(h)-\min\{1,\epsilon\}^2/288. \tag{3.3}
\]

The margin remains positive at h=1/4. It is a new geometric scope for
the energy wall, not a claim that every topology forces its candidates
into the excluded band. In particular, (3.2) does not supply an Einstein
existence theorem from failure of the obstruction.

\hypertarget{euler-density-for-curvature-regular-elliptic-metrics}{%
\subsection{4. Euler density for curvature-regular elliptic
metrics}\label{euler-density-for-curvature-regular-elliptic-metrics}}

One can already remove the explicit covering premise in a distinct,
classical elliptic class.
\href{https://nyjm.albany.edu/j/2017/23-64p.pdf}{Weber, Theorem 1.1}
gives collapse-independent, unnormalized small-energy regularity for the
critical four-metrics specified there, including Einstein metrics,
constant-scalar-curvature Bach-flat metrics and metrics with harmonic
curvature. Here use these explicitly covered classes; ``critical'' is
not defined to include an arbitrary metric or an arbitrary variational
PDE.

Write R\_op for the curvature operator, P for integral
\textbar Ric\^{}circ\textbar\^{}2, and Lambda=8 pi\^{}2 chi(M)/V.
Four-dimensional Chern--Gauss--Bonnet gives

\[
                  \int|R_{\rm op}|^2=8\pi^2\chi+P.   \tag{4.1}
\]

For such metrics with \texttt{\textbar{}Ric\textbar{}\textless{}=K} and
\texttt{Lambda\textless{}=Lambda\_0}, the normalized total curvature
energy is at most Lambda\_0+K\^{}2. Fubini and the Bishop ball-volume
upper bound show that, at a fixed sufficiently small radius
R=R(K,Lambda\_0), at least half the centers have curvature energy below
Weber's threshold. Its conclusion bounds sectional curvature on a
uniformly smaller ball. The exponential-development averaging theorem of
the preceding manuscript, Section 11, then gives

\[
                   \fint Z_t\ge c(K,L,\Lambda_0)>0.  \tag{4.2}
\]

The proof is exactly the curvature-energy-to-regular-fraction argument
in Section 12 there, with (4.1) replacing the exact-Einstein identity
and Weber replacing Cheeger--Tian. It does not require that the Ricci
error vanish identically. Composing with (2.1) gives the corresponding
actual gap. This is a genuine approximate-input extension, but keeps the
stated elliptic metric restriction. The next section gives a verified
smoothing route beyond that restriction, with a different, uniform local
energy input. It does not replace that input by an averaged
Euler-density bound.

\hypertarget{smoothing-supplies-the-covering-geometry-without-moving-the-form}{%
\subsubsection{4.1. Smoothing supplies the covering geometry, without
moving the
form}\label{smoothing-supplies-the-covering-geometry-without-moving-the-form}}

\textbf{Comparison-metric lemma.} Suppose the original metric g has
bounded Ricci and admits a complete metric h on the same manifold such
that

\[
                      a^{-1}h\le g\le ah,
                  \qquad |\sec_h|\le H.              \tag{4.3}
\]

Then there are rho=rho(a,H)\textgreater0 and v=v(a)\textgreater0 for
which EVERY center satisfies w\_rho\textgreater=v in (1.2). In
particular

\[
                         \fint Z_t\ge c(K,L,a,H)>0.   \tag{4.4}
\]

\textbf{Proof.} Use
\href{https://www.ihes.fr/~gromov/wp-content/uploads/2018/08/384.pdf}{Cheeger--Fukaya--Gromov,
Theorem 1.3 and definition (1.1.4)--(1.1.5)}, printed pages 331--332, on
the complete metric Sh, S=max(1,2H), with approximation parameter log 2.
After rescaling it supplies a metric k, uniformly equivalent to h within
factor 2 and with \textbar sec\_k\textbar\textless=C\_0 S. At every x
there is an actual normal cover of an open set V containing
B\_s\^{}k(x), s=rho\_0/sqrt(S), on which the lifted k-metric has
injectivity radius at least s. This is the covering assertion in the
definition of their round metric, justified on printed page 365; it is
not a claim that an exponential chart is a cover.

Write a'=2a, so
\texttt{(a\textquotesingle{})\^{}-1\ k\textless{}=g\textless{}=a\textquotesingle{}\ k}.
Choose a dimensional d\_0 below rho\_0/2 and 1/(2 sqrt(C\_0)), and let
\texttt{rho=d\_0/(8\ sqrt(a\textquotesingle{}\ S))}. Then
\texttt{U=B\^{}g\_\{4rho\}(x)} lies inside B\_s\^{}k(x). The universal
cover of U maps, by the lifted inclusion, to the supplied normal cover
of V. This map F is a local isometry for the ORIGINAL g. Its image on
B\_rho contains the g-radius-rho ball of the normal cover: every path of
g-length less than rho starting at the marked lift projects into
B\_rho\^{}g(x), and thus lifts again to the universal cover of U.
Uniqueness of path lifting gives the asserted surjection on the balls.

The normal cover's g-ball contains its k-ball of radius rho/sqrt(a').
This radius is below its injectivity and curvature scales. Rauch
comparison and
\texttt{dV\_g\textgreater{}=(a\textquotesingle{})\^{}-2\ dV\_k} give
volume at least \texttt{c\_4\ (a\textquotesingle{})\^{}-4\ rho\^{}4}.
The area formula for the surjective local isometry gives the same lower
bound on the universal cover of U. All these balls have an interior
buffer: a finite-length escaping path cannot lie over an interior limit
point of a covering base, by an evenly covered coordinate neighborhood.
Now apply Theorem 1 with G=M. QED.

The metric h is used only to supply covering volume. The supplied form,
its self-duality, the Ricci error, and the curvature deficit are all
evaluated in g. No assertion that these errors are preserved by
smoothing is needed.

\textbf{Fixed-scale nonconcentration corollary.} There is a dimensional
epsilon\_L\textgreater0 such that, for every fixed
0\textless r\textless=1, actual metrics satisfying delta\_-\textless=1/4
and

\[
             \sup_x\int_{B_r^g(x)}|\operatorname{Rm}_g|^2
                              \le\epsilon_L          \tag{4.5}
\]

obey

\[
 \boxed{\delta_-^2+D_\beta/V\ge c(r)>0,
                              \qquad\beta\ge4.}      \tag{4.6}
\]

With the budgets of Section 3 this also gives
\texttt{delta\_-\^{}2+(G(h)-e)\_+\textgreater{}=c\textquotesingle{}(r)\textgreater{}0}
on \texttt{1/4\textless{}=h\textless{}=1/3}. No diameter, noncollapse,
Euler-density, class-square, exact Einstein, or critical-metric
assumption is added. Here Rm is the FULL four-index tensor:
\texttt{\textbar{}Rm\textbar{}\^{}2=4\textbar{}R\_op\textbar{}\^{}2};
adjust epsilon\_L by four when using the curvature-operator convention.

To prove this, \href{https://arxiv.org/pdf/0911.3104}{Li II}, Theorem
1.1 and Section 3, provides a definite-time smoothing under (4.5) and a
Ricci bound. Its metric comparison (3.8) is closed in the proof, and at
a fixed positive time it gives (4.3) with a=2 and
\texttt{H\textless{}=C\ max(r\^{}-2,K)}. Compose the comparison lemma
with the actual supplier (2.1), and then with Section 3's algebra.

For source precision: the available arXiv v1 states the complete
noncompact theorem, whereas the
\href{https://link.springer.com/article/10.1007/s10455-011-9290-0}{published
primary abstract} explicitly covers compact manifolds. The published
full text was not retrieved. The compact application also follows from
the displayed Section 3 bootstrap with Hamilton's closed initial
existence and extension: when delta\_-\textless=1/4,
Chern--Gauss--Bonnet has
\texttt{s\^{}2/24-\textbar{}Ric\^{}circ\textbar{}\^{}2/2\textgreater{}=1/4},
hence chi\textgreater0 and
\texttt{integral\ \textbar{}Rm\textbar{}\^{}2\textgreater{}=32\ pi\^{}2}.
Decrease epsilon\_L below this. Condition (4.5) forces
diam(M)\textgreater r; run the local Dirichlet-Sobolev bootstrap on
balls at scale r/10, all properly buffered. The finite-radius covering
and energy estimates then apply unchanged. The integrable Ricci bound
\texttt{C\ t\^{}(-2/3)} in that proof gives the factor-two metric
comparison by integrating in time. Only dimensional constants are
reduced. Thus the use of the compact case is explained, not silently
quoted from a noncompact-only statement.

A directly checkable sufficient hypothesis is a local supercritical
curvature budget: for some eta\textgreater0,

\[
                \sup_x\int_{B_1(x)}|\operatorname{Rm}|^{2+\eta}
                           \le\Lambda.               \tag{4.7}
\]

Indeed Bishop comparison and Holder give, uniformly at every center,

\[
 \int_{B_r(x)}|\operatorname{Rm}|^2
 \le\Lambda^{2/(2+\eta)}
                [C(K)r^4]^{\eta/(2+\eta)}.           \tag{4.8}
\]

Choose a fixed r making this less than epsilon\_L. Thus (4.6) and the
upper wall hold with constants depending on eta,Lambda, without a
pointwise full-curvature bound. Any specified uniform local L2
equiintegrability modulus also suffices.

Conversely, a pointwise near-Einstein actual zero-end sequence must lose
(4.5) at every fixed radius eventually. Choosing radii tending to zero
diagonally gives actual balls carrying at least epsilon\_L of full
curvature energy. This is a necessary concentration statement, not a
classification or construction of its possible bubbles.

An averaged bound chi/V controls total L2 energy per volume, not the
supremum over centers in (4.5). Without an elliptic metric equation,
extending Section 4's Euler-density conclusion requires a verified
buffered LOCAL smoothing and covering argument. The abstract of Li I
alone is not used to declare that stronger extension proved.

\hypertarget{exact-all-metric-defect-on-the-handle-only-family}{%
\subsection{5. Exact all-metric defect on the handle-only
family}\label{exact-all-metric-defect-on-the-handle-only-family}}

Let X be a closed complex-hyperbolic surface in its complex orientation,
\texttt{A=c1\^{}2(X)=3\ chi(X)\textgreater{}0}. Work in the established
actual-monopole surgery class

\[
 M=X\#k\overline{\mathbb {CP}}^{\,2}
           \#\ell(S^1\times S^3),\qquad
 h=(k+4\ell)/A,\qquad
 \theta=h-6k/A=(4\ell-5k)/A.                       \tag{5.1}
\]

The metricwise scalar estimate supplies S\textgreater=S\_0=32 pi\^{}2 A,
where \texttt{S=integral\ s\^{}2}. Define

\[
 B=\int|W^-|^2,\quad P=\int|\operatorname{Ric}^{\circ}|^2,
 \quad L_s=\int(s-\bar s)^2,\quad R=P+S/4.
\]

These assumptions and the classical gluing construction are the same as
in the preserved complex-hyperbolic application; neither formal classes
nor formal curvature tensors are substituted for actual metrics.

\textbf{Proposition 3 (signed scalar--Weyl identity).} If
0\textless=theta\textless=1/3, then for EVERY smooth metric on M,

\[
 \boxed{\mathcal E-\frac{4\theta}{1+\theta}
   =\frac{(1-3\theta)(S-S_0)/3+16B+(1+\theta)L_s}
                {(1+\theta)R}.}                     \tag{5.2}
\]

In particular the elementary all-metric lower bound is

\[
 \mathcal E\ge
 \begin{cases}
 0,&\theta\le0,\\
 4\theta/(1+\theta),&0\le\theta\le1/3,\\
 1,&\theta\ge1/3.
 \end{cases}                                         \tag{5.3}
\]

\textbf{Proof.} Signature and the positive Chern--Weil identity give

\[
 \int|W^+|^2=B+12\pi^2(A/9-k),\qquad
 P=4\int|W^+|^2+S/12-8\pi^2(A-k-4\ell).
\]

Hence

\[
 P=S/12+4B+8\pi^2A(\theta-1/3),\qquad
 \mathcal E=(4P+L_s)/R.                              \tag{5.4}
\]

Substitute (5.4) in (5.2). Every numerator term is nonnegative on the
stated theta interval. When theta\textgreater=1/3 instead use

\[
 \mathcal E-1=
 \frac{8\pi^2A(3\theta-1)+12B+L_s}{R}\ge0.           \tag{5.5}
\]

R is nonzero by the positive scalar budget. QED.

\textbf{Corollary 4 (exact infimum, no restriction to constructed
metrics).} For k=0 and \texttt{0\textless{}=h=4ell/A\textless{}=1/3},

\[
                  \boxed{\inf_g\mathcal E(g)=\frac{4h}{1+h}.} \tag{5.6}
\]

The lower bound follows from (5.2), now theta=h. For the upper bound,
the classical shrinking-core gluing sequence on this SAME smooth
manifold satisfies

\[
 R\to8\pi^2A(1+h),\quad
 V\to2\pi^2A,\quad \int s\to-8\pi^2A,
\]

so \texttt{E=4-(integral\ s)\^{}2/(VR)} tends to the right side of
(5.6). The use of this classical construction, originating in
\href{https://arxiv.org/abs/math/0003068}{LeBrun's Ricci-energy
calculation}, is already detailed in the boundary continuation. This is
an exact calculation of the project functional from classical budgets
and construction, not a new gluing theory or another proof of I40.

If h=0 the complex-hyperbolic metric attains the infimum. If
0\textless h\textless1/3, equality in (5.2) forces S=S\_0, B=0 and
L\_s=0. The actual ordinary monopole scalar equality then forces a
nonzero parallel positive spinor. For the handle class its square is
a\^{}2=A and

\[
                         c_2(S^+)[M]=(a^2-T_+)/4=\ell>0,
\]

contradicting a nowhere-zero section. Thus the infimum in this open
interval is not attained. This argument does NOT infer Einstein from the
cscK equality statement. At h=1/3 the first coefficient in (5.2)
vanishes; the exact identity is \texttt{E-1=(12B+L\_s)/R}, and
attainment is not decided by the preceding argument.

More generally, along surgery families with 00 and \texttt{k=o(Ah)},
(5.3) and the same classical upper construction yield

\[
                         \inf_g\mathcal E(g)/h\longrightarrow4. \tag{5.7}
\]

For the balanced tight-window family, however,
\texttt{k/A=2h/9+O(h\^{}2)}, so theta=-h/3+O(h\^{}2). Bound (5.3) is
then zero. The all-metric linear gap proved by I40--I41 is not replaced
by this classical calculation; its best coefficient remains genuinely
open on that family.

\hypertarget{a-chiral-blowdown-cap-theorem-with-no-scalar-cost-on-the-caps}{%
\subsection{6. A chiral blowdown-cap theorem, with no scalar cost on the
caps}\label{a-chiral-blowdown-cap-theorem-with-no-scalar-cost-on-the-caps}}

For an arbitrary metric on M, suppose only the k blowup cores are
removed and replaced by compatible smooth caps, giving

\[
                         N=X\#\ell(S^1\times S^3).
\]

Require a smooth metric g\_N equal to g\_M on the full common exterior,
including matched collars. All handle geometry is left unchanged. A
smooth extension is not itself an energy bound. Put

\[
 B_N=\int_N|W_N^-|^2,\quad
 K_+=\int_{\mathrm{cores}_M}|W_M^+|^2,\quad
 C_+=\int_{\mathrm{caps}_N}|W_N^+|^2.                 \tag{6.1}
\]

\textbf{Theorem 5 (Weyl-only cap cost).} For 0\textless=h\textless=1/3,
the exact identity is

\[
 \boxed{\mathcal E_M-\frac{4h}{1+h}
  =\frac{(1-3h)(S_M-S_0)/3+16(B_N+K_+-C_+)+(1+h)L_M}
                        {(1+h)R_M}.}                 \tag{6.2}
\]

In particular, under the inherited actual scalar/projection budget,

\[
 \boxed{\mathcal E_M\ge\frac{4h}{1+h}
       -\frac{16C_+}{(1+h)R_M}
       \ge\frac{4h}{1+h}-\frac{2C_+}{\pi^2 A(1+h)^2}.} \tag{6.3}
\]

Neither scalar, Ricci, opposite Weyl nor volume bounds on the caps are
required. The cap metric does not have to solve a gauge or Einstein
equation. The proof uses actual scalar information on M, not on N.

\textbf{Proof.} \texttt{tau(N)-tau(M)=k}. Subtract the signature formula
for the two closed metrics, cancelling their common exterior. This gives

\[
                 B_M=12\pi^2k+B_N+K_+-C_+.           \tag{6.4}
\]

Substitute this in (5.4), noting theta=h-6k/A. The k terms cancel:

\[
       P_M=8\pi^2Ah+(S_M-S_0)/12+4(B_N+K_+-C_+).
                                                        \tag{6.5}
\]

Using \texttt{E=(4P+L)/R} proves (6.2). All its numerator terms except
-16C\_+ are nonnegative. Finally the classical beta=2 bound gives
\texttt{R\_M\textgreater{}=8\ pi\^{}2\ A(1+h)} (equivalently
\texttt{D\_2=36\ pi\^{}2\ A(e-1-h-2z)\textgreater{}=0}, where
\texttt{z=p/A-1\textgreater{}=0}). This proves the second inequality in
(6.3). QED.

Along 00, \texttt{C\_+=o(Ah)} suffices for
\texttt{liminf\ E/h\textgreater{}=4}, without Yamabe reduction or the
tight opposite-chirality window. If compatible caps are ASD, then C\_+=0
and the finite bound is \texttt{E\textgreater{}=4h/(1+h)}. This is a
conditional metric theorem: existence of ASD or arbitrarily cheap
compatible caps for all prescribed boundary metrics is not asserted.

The weaker signed condition \texttt{C\_+\textless{}=B\_N+K\_+} also
suffices for the same finite lower bound. In particular nonzero cap
energy is allowed; only its uncompensated excess matters.
Four-dimensional chiral Weyl energy is conformally invariant. A
conformal change of a filling can repair the conformal factor of its
matching collar, but cannot reduce C\_+.

For comparison, subtracting positive Chern--Weil instead of signature
would give only

\[
 P_M\ge8\pi^2(k+4\ell)
          -2\int_{\mathrm{caps}_N}(s_N^2/24+2|W_N^+|^2).
\]

Retaining the signature relation removes the entire scalar cap cost.
This is why controlling all curvature of a cap, or also reversing the
handles, is unnecessarily strong for this target.

Two concrete sufficient regimes for (6.3) are:

\begin{itemize}
\tightlist
\item
  Compatible caps with
  \texttt{sum\ Vol(C\_i)\textbar{}\textbar{}W\^{}+\textbar{}\textbar{}\_infinity,C\_i\^{}2=o(Ah)}.
\item
  Small almost-Euclidean boundary collars, in scaled coordinates
  \texttt{ell\_i\^{}2(delta+q\_i)} with
  \texttt{\textbar{}\textbar{}q\_i\textbar{}\textbar{}C\^{}2\textless{}=epsilon\_i},
  capped by a smooth cutoff of q\_i. Then
  \texttt{C\_+\textless{}=C\ sum\ epsilon\_i\^{}2}; requiring this sum
  to be o(Ah) is enough. For uniformly smooth ambient collars
  \texttt{epsilon\_i=O(ell\_i\^{}2)}, \texttt{sum\ ell\_i\^{}4=o(Ah)}
  suffices. Since k\textless=Ah, a uniform
  \texttt{max\ ell\_i\ -\textgreater{}\ 0} implies that condition.
\end{itemize}

The cutoff statement follows by scaling to the unit ball, extending the
given small collar tensor smoothly to a slightly larger annulus with
uniform C2 control, and setting it to zero in the interior. Curvature is
O(epsilon\_i), its squared integral is O(epsilon\_i\^{}2), and the
four-dimensional Weyl integral is scale invariant. This construction
requires the stated controlled collars; small total Ricci error alone
does not supply them.

Conversely, along ANY such actual surgery family with 00, if
\texttt{E/h-\textgreater{}c\textless{}4}, the second inequality of (6.3)
forces EVERY compatible blowdown cap system to satisfy

\[
                \liminf\frac{C_+}{Ah}
                              \ge\frac{\pi^2}{2}(4-c). \tag{6.6}
\]

This is a necessary chiral boundary cost for any improvement below 4,
not a proof that the cost must vanish. No scalar normalization or limit
\texttt{R/(8\ pi\^{}2\ A)-\textgreater{}1} is needed for this
conclusion.

Equivalently, for any scalar -4 candidate, (5.4) isolates the exact
remaining quantity:

\[
 P-8\pi^2Ah
 =4\left\{\frac{V-2\pi^2 A}{3}
                    +\int|W^-|^2-12\pi^2 k\right\}. \tag{6.7}
\]

Thus the optimal coefficient question asks whether the volume excess and
opposite Weyl energy must pay the blowup cost, up to o(Ah). It is not a
request to eliminate every ordinary-spinor remainder separately, nor to
cap every surgery in the construction.

\hypertarget{exact-cost-of-a-prescribed-conformal-complex-hyperbolic-collar}{%
\subsubsection{6.1. Exact cost of a prescribed conformal
complex-hyperbolic
collar}\label{exact-cost-of-a-prescribed-conformal-complex-hyperbolic-collar}}

The sufficient filling condition is geometric, not merely topological.
Its cost can be computed exactly for an important full collar germ. Let
a marked cap be a four-ball with prescribed metric on a full boundary
collar, conformal there to the germ of the radius-r ball in complex
hyperbolic two-space with Ric=-g. Preserve the ball's complex
orientation and its boundary marking. Among all smooth positive-definite
extensions of this exact collar metric,

\[
 \boxed{\inf_{g_C}\int_C|W_{g_C}^+|^2
              =\frac23\operatorname{Vol}_{\mathbb{CH}^2}(B_r).} \tag{6.8}
\]

\textbf{Proof.} Extend the prescribed positive conformal factor smoothly
over the ball, producing a reference extension g\_0 conformal to the
complex hyperbolic metric. There is no derivative bound on this factor.
Since \texttt{W\^{}-=0} and
\texttt{\textbar{}W\^{}+\textbar{}\^{}2=s\^{}2/24=2/3} for the reference
metric before conformal change, conformal invariance gives its chiral
energy in (6.8).

Define
\texttt{b(g)=integral\_C(\textbar{}W\^{}+\textbar{}\^{}2-\textbar{}W\^{}-\textbar{}\^{}2)=4\ pi\^{}2\ integral\_C\ p\_1(nabla\^{}g)}.
For any other extension g\_C, canonically identify the two orthonormal
frame bundles; the identification is the identity on their common full
collar. The relative Pontryagin transgression has zero boundary value,
so \texttt{b(g\_C)=b(g\_0)}. Consequently

\[
 \int_C|W_{g_C}^+|^2
       =\frac23\operatorname{Vol}_{\mathbb{CH}^2}(B_r)
                      +\int_C|W_{g_C}^-|^2.          \tag{6.9}
\]

The reference extension attains the lower bound. This proof compares
connections on the SAME marked ball; it does not double two inward
collars without checking their normal derivatives. QED.

For several such collars the exact minimum is the sum of (6.8). Thus a
total reference ball volume o(Ah) is sufficient and necessary for cheap
caps of this specified conformal type. Fixed positive radii on a
positive proportion of k comparable to Ah force order-Ah cost for these
prescribed cuts, even if the widths of the prescribed collars tend to
zero. For full collars conformal to a flat round-ball germ, the same
construction gives an exactly conformally flat cap and C\_+=0. An
intrinsically round boundary metric alone is not this full-germ
hypothesis: the normal derivatives of the metric also enter the relative
Pontryagin boundary data.

These statements do not assert that alternative cuts through a given
exceptional neighborhood have the same cost. Finding compatible cuts
with sufficiently small uncompensated chiral cost remains part of the
all-metric coefficient problem.

\hypertarget{round-induced-boundary-does-not-imply-a-cheap-compatible-cap}{%
\subsubsection{6.2. Round induced boundary does not imply a cheap
compatible
cap}\label{round-induced-boundary-does-not-imply-a-cheap-compatible-cap}}

There are smooth metrics on the standard four-ball with the SAME unit
round induced boundary metric for which every full-collar-compatible cap
has arbitrarily large positive chiral energy. Thus the full-germ clause
in Section 6.1 cannot be replaced by intrinsic boundary geometry alone.

Take a coframe on the unit three-sphere with
\texttt{d\ sigma\_i=-2\ sigma\_j\ wedge\ sigma\_k} cyclically and
\texttt{integral\ sigma\_1\ sigma\_2\ sigma\_3=2\ pi\^{}2}. For
t\textgreater0 put

\[
 g_t=dr^2+r^2(\sigma_1^2+\sigma_2^2)
        +r^2F(r)^2\sigma_3^2,
 \qquad F(r)=1+t r^2(1-r^2),\quad 0\le r\le1.       \tag{6.10}
\]

This is smooth at the origin: its difference from the Euclidean metric
is a smooth radial polynomial times the square of the smooth one-form
\texttt{r\^{}2\ sigma\_3}. It is positive definite. At r=1 its induced
metric is round, but its outward principal curvatures are
\texttt{(1,1,1-2t)}.

Here is an explicit normalization check. With orientation
\texttt{dr\ wedge\ sigma\_1\ wedge\ sigma\_2\ wedge\ sigma\_3}, write
G=F+rF'. In the orthonormal coframe
\texttt{(dr,r\ sigma\_1,r\ sigma\_2,rF\ sigma\_3)}, the connection is

\[
 \omega_{10}=\sigma_1,\quad\omega_{20}=\sigma_2,
 \quad\omega_{30}=G\sigma_3,\quad
 \omega_{12}=(F^2-2)\sigma_3,\quad
 \omega_{13}=F\sigma_2,\quad\omega_{23}=-F\sigma_1.
\]

Computing \texttt{Omega=d\ omega+omega\ wedge\ omega} gives the signed
density

\[
 \begin{aligned}
 (|W^+|^2-|W^-|^2)d\mu
 &= -\tfrac12\operatorname{tr}(\Omega\wedge\Omega)\\
 &=\frac{d}{dr}\{-4(1-F^2)^2-2r^2(F')^2\}
                       \,dr\wedge\sigma_1\wedge\sigma_2\wedge\sigma_3.
 \end{aligned}\tag{6.11}
\]

Since F(0)=F(1)=1 and F'(1)=-2t, its integral is
\texttt{-16\ pi\^{}2\ t\^{}2}. Pull the metric back by a reflection of
the ball, retaining the original orientation of the domain. The
reflected metric still induces the round boundary metric and has signed
integral \texttt{+16\ pi\^{}2\ t\^{}2}.

Relative transgression on the marked ball now proves that every metric
matching this reflected FULL collar satisfies

\[
                     \int_C|W^+|^2\ge16\pi^2t^2.    \tag{6.12}
\]

The local obstruction is exact and survives overall rescaling and
conformal changes. It does not obstruct choosing a different cut farther
out, and does not give a metric with E/h\textless4. In particular it is
not a counterexample to the unresolved all-metric coefficient
conjecture.

\hypertarget{global-first-order-ale-matching-a-finite-energy-obstruction}{%
\subsection{7. Global first-order ALE matching: a finite-energy
obstruction}\label{global-first-order-ale-matching-a-finite-energy-obstruction}}

The positive-charge local cores in the preceding continuation solve the
actual weighted equations, but are not global solutions. Here we analyze
a precise global matching limit, retaining its first-order connection
source rather than declaring a twisted Dirac kernel trivial.

\textbf{Theorem 6 (first-order ALE matching).} Let (X,h) be a complete
one-ended hyperkaehler ALE four-manifold, with W\_h\^{}+=0, nontrivial
free end group \texttt{Gamma\ subset\ SU(2)\_-}, and
\texttt{H\^{}1(X;R)=0}. Assume the FULL limiting positive Spin-c
connection, including its determinant twist, admits a global parallel
unitary frame. On an exhaustion suppose actual rescaled fields have

\[
 \epsilon_j^{-1}\operatorname{Ric}^{\circ}(h_j)\longrightarrow0,\quad
 h_j\longrightarrow h,\quad
 D_{h_j,A_j}\Phi_j=0,\quad
 F_{A_j}^{+,h_j}=i\epsilon_j^2 f_j\sigma_{h_j}(\Phi_j),
 \qquad \epsilon_j\downarrow0,                       \tag{7.1}
\]

and compatible global bundle identifications give

\[
 \nabla_j^+=d+\epsilon_jB+o(\epsilon_j),\qquad
 \Phi_j=\Phi_0+\epsilon_j\psi+o(\epsilon_j),\qquad
 \epsilon_j f_j\longrightarrow0.                    \tag{7.2}
\]

All displayed convergence is local, with the weight convergence locally
uniform. The metric convergence is smooth, the connection expansion is
C1, and the spinor expansion is C2. Phi\_0 is nonzero and parallel. B is
a global u(2)-valued one-form in the given frame. Require only

\[
             dB\in L^2(X),\qquad
             U=d\psi+B\Phi_0\text{ bounded on the ALE end}.   \tag{7.3}
\]

Then U=0. In particular the limiting physical cubic torsion

\[
             T=\sum_i c_h(e^i)\Phi_0
                             \otimes(\Phi_0\wedge U_i)        \tag{7.4}
\]

vanishes. Nonzero calibrated exterior torsion cannot occur in this same
global first-order limit. No pointwise decay of B, or growth bound on
the ungauged spinor correction psi, is required.

\textbf{Proof.} Split B into real components b\_alpha in a constant u(2)
basis, with b\_0 central. The background connection is flat, so its
first-order curvature is dB. The first hypothesis in (7.1) makes the
FIRST-ORDER trace-free positive-spinor curvature self-dual as a
two-form, while the scaled weighted equation gives anti-self-dual
first-order central curvature:

\[
            (db_\alpha)^-=0\quad(\alpha=1,2,3),
                    \qquad(db_0)^+=0.               \tag{7.5}
\]

We need an ALE primitive fact: a smooth globally exact, L2, self-dual or
anti-self-dual two-form F=db on X vanishes. Lift the end to an exterior
domain in R4 and extend F to R4 as \texttt{F\_ext=d(chi\ b)}, where chi
vanishes on a fixed inner neighborhood and equals one far out. This
extension is closed and L2. Euclidean Fourier/Hodge inversion gives

\[
        \beta=d^*\Delta^{-1}F_{\rm ext},\qquad
        d\beta=F_{\rm ext},\quad
        \|\nabla\beta\|_2\le C\|F_{\rm ext}\|_2,
        \quad\beta\in L^4(\mathbb R^4).              \tag{7.6}
\]

The last assertion is homogeneous Euclidean Sobolev, not a uniform
Sobolev premise on a degenerating closed manifold. Average over the
finite end group to obtain a descending primitive. On the end b-beta is
closed and, since \texttt{H\^{}1(S\^{}3/Gamma;R)=0}, equals dq.
Subtracting d(eta q), where eta is supported on the end and equals one
far out, produces a global primitive b\_tilde in L4 with db\_tilde=F.
Euclidean and ALE norms are equivalent there. For radial cutoffs chi\_R,
the L4 norm of dchi\_R is uniformly bounded, so Stokes and Holder yield

\[
 \left|\int_X\chi_R F\wedge F\right|
 \le C\|\tilde b\|_{L^4(\{R<r<2R\})}
         \|F\|_{L^2(\{R<r<2R\})}\longrightarrow0.   \tag{7.7}
\]

Definite chirality identifies the limit with plus or minus the total L2
energy. Thus F=0. Apply this fact componentwise to (7.5), obtaining
dB=0. Since H1(X;R)=0, B=dG for a global skew-Hermitian matrix function;
G need not decay or be bounded.

The first-order Dirac equation is
\texttt{D\_h\ psi=-sum\_i\ c\_h(e\^{}i)\ B\_i\ Phi\_0}. Variation of the
principal symbol multiplies dPhi\_0=0; all connection variations have
been retained in B. Therefore

\[
            \psi'=\psi+G\Phi_0,\qquad D_h\psi'=0,
                                    \quad d\psi'=U.  \tag{7.8}
\]

The bounded derivative makes psi' grow at most linearly. Lichnerowicz in
the global parallel frame says its components are scalar harmonic
functions. Every such function with at most linear growth is constant on
this ALE space: blowdowns u(Rx)/R have subsequential harmonic limits on
the Euclidean end cover bounded by C\textbar x\textbar. The origin is
removable; an entire harmonic function of linear growth is affine. The
limit vanishes at the origin and is Gamma-invariant, hence is linear and
Gamma-invariant. A nontrivial group acting freely on S3 has no fixed
vector, so the limit is zero. Every blowdown vanishes, giving u=o(r).
The Ricci-flat interior gradient estimate then gives du=0. Consequently
dpsi'=U=0. QED.

Bounded U can be replaced by bounded limiting physical torsion T and
current \texttt{Z=\textless{}Phi\_0,U\textgreater{}}. The linearized
exact Dirac decomposition, with rho\_0=\textbar Phi\_0\textbar\^{}2,
gives

\[
        0\le2\rho_0^2|U|^2-|T|^2
                         \le4\rho_0|Z|^2.            \tag{7.9}
\]

In particular an exterior with
\texttt{\textbar{}T\textbar{}\^{}2\ -\textgreater{}\ 128/81} and
Z-\textgreater0 cannot be the exterior of this same global first-order
limit. Exact Einstein is NOT required: for g\_j=epsilon\_j\^{}2 h\_j,
covariant Ricci and its trace-free part are unchanged, but

\[
 |\operatorname{Ric}^{\circ}(h_j)|_{h_j}
       =\epsilon_j^2|\operatorname{Ric}^{\circ}(g_j)|_{g_j}.
                                                        \tag{7.10}
\]

Thus uniformly bounded PHYSICAL Ricci supplies the first condition in
(7.1). It also supplies the canonical weight condition. For the
inherited majorant its exact rescaling is

\[
 \epsilon_j f_{0,j}=\max\left\{
 \sqrt{\max\{-s_{h_j}+\sqrt6|W^+_{h_j}|,3\epsilon_j^2\}},
                 \frac{(-s_{h_j})_+}{4\epsilon_j}\right\}\longrightarrow0,
                                                        \tag{7.10a}
\]

since s\_h=O(epsilon\^{}2) and W\_h\^{}+ tends to zero locally. Close
smooth majorants preserve this property. With only the trace-free
hypothesis in (7.1), the separate weight limit in (7.2) must be
retained: it does not follow just by ignoring the majorant's scalar
branch.

The stronger pointwise conditions
\texttt{B=O(r\^{}-1-delta),\ dB=O(r\^{}-2-delta),\ psi=O(r)} give a
simpler proof by boundary Stokes and a decaying gauge primitive, but are
unnecessary in the finite-energy version. For the flat complete limit R4
the harmonic components can be affine, so U and T are constant. A
positive simple torsion zero still cannot persist in a bounded global
first-order limit. The earlier local core has linearly growing rescaled
torsion, not this bounded exterior matching.

The original averaged budgets do NOT yet supply first-order compactness
at scale epsilon, finite total first-order curvature energy, or nonzero
exterior torsion visible in that same limit. Physical curvature of size
epsilon\^{}-1 on an epsilon-ball has only O(epsilon\^{}2) ordinary L2
cost; that does not bound first-order curvature on an entire expanding
end. Calibration may move to radii R\_j-\textgreater infinity in bubble
coordinates. These limits cannot be interchanged.

The full trivial limiting Spin-c connection is essential. A parallel
frame of the untwisted spin bundle with an unspecified flat determinant
character would not make the correction components Gamma-invariant
scalar functions.

\hypertarget{bare-signed-degree-is-not-the-missing-global-constraint}{%
\subsubsection{7.1. Bare signed degree is not the missing global
constraint}\label{bare-signed-degree-is-not-the-missing-global-constraint}}

For \texttt{E\_-=S\^{}-\ tensor\ L}, the inherited identity is
\texttt{c\_2(E\_-){[}M{]}=(9a\^{}2+T\_-)/4}. If the exterior section is
nonzero and N balls contain positive simple torsion zeros, necessarily

\[
                     N=(9a^2+T_-)/4.                 \tag{7.11}
\]

At the zero endpoint this becomes

\[
  N/V=\frac1{8\pi^2V}\int|W^-|^2
                            +\frac5{12\pi^2}+o(1),   \tag{7.12}
\]

which has the sign compatible with positive cores. Conversely, as a
SECTION problem, the relative Euler obstruction is the only one: the
unit-sphere fiber S3 is 2-connected and the exterior has dimension four.
When (7.11) holds, prescribed nonzero boundary germs with these degrees
extend to a nonzero exterior section. This does not make the extension
the cubic torsion of a monopole. It explains why Theorem 6 retains the
actual Dirac and curvature equations rather than only zero counts.

\hypertarget{what-these-connections-decide-and-what-they-do-not}{%
\subsection{8. What these connections decide, and what they do
not}\label{what-these-connections-decide-and-what-they-do-not}}

The original obstruction mechanism is unchanged: actual fields detect
curvature modes incompatible with the differential equations. This
continuation identifies additional geometric situations in which that
detection remains quantitative, computes an exact value on a genuine
subfamily, and determines which filling/matching information the next
sharpness and endpoint arguments need.

The general zero-end sequence would have to escape both local-cover
control and the verified smoothing classes. The balanced family's
coefficient below 4 would have to pay a definite uncompensated chiral
filling cost for every compatible blowdown. A local positive torsion
core cannot be matched through the finite-energy, bounded-derivative ALE
first-order model of Section 7. None of these necessary conditions is
itself a constructed escaping sequence.

Still unresolved are the unrestricted weak-error zero endpoint, the
balanced family's all-metric optimal coefficient, and the global
Einstein decision on fixed E\_D. The proof does not reclassify these as
completed merely because their remaining mechanisms are now more
specific. The new exact subfamily formula and transgression computations
retain their classical-input status; worldwide novelty is a separate
literature task.

END COMPLETE SOURCE

\begin{center}\rule{0.5\linewidth}{0.5pt}\end{center}

\hypertarget{source-2-20260927_local_geometry_and_variational_extrema.md}{%
\subsection{Source 2:
20260927\_local\_geometry\_and\_variational\_extrema.md}\label{source-2-20260927_local_geometry_and_variational_extrema.md}}

Original path:
continuation/20260927\_local\_geometry\_and\_variational\_extrema.md

SHA-256:
68abdebe39870e7e984169eb9b434a88b0a974b98da8e10b7effa494e7c1b7e6

BEGIN COMPLETE SOURCE

\hypertarget{local-unwrapping-annular-matching-and-variational-extrema}{%
\section{Local unwrapping, annular matching, and variational
extrema}\label{local-unwrapping-annular-matching-and-variational-extrema}}

27 September 2026. Continuation from \texttt{cee6c6b}, on the local
branch \texttt{local/yau26-variational-global-20260927}. This file
extends the
\href{20260927_covering_geometry_and_sharp_subfamilies.md}{covering and
sharp-subfamily proof}. It does not replace the frozen manuscripts or
create a web-export packet.

The inherited linear small-h order and scalar-variance removal are
premises, not new results here. Curvature norms, actual-class
quantifiers and positive weight solvability have the conventions of the
\href{20260927_conformal_separation_and_zero_end.md}{conformal
continuation}. Statements about smooth minimizers below are conditional
on attainment; they do not supply compactness or a minimizer.

\hypertarget{the-exact-variational-problem-behind-the-remaining-coefficient}{%
\subsection{1. The exact variational problem behind the remaining
coefficient}\label{the-exact-variational-problem-behind-the-remaining-coefficient}}

This section gives a usable connection to the classical
quadratic-curvature problem, rather than a new lower bound. Write

\[
 V=\operatorname{Vol}(g),\quad S_1=\int s,\quad S_2=\int s^2,
 \quad R=\int|\operatorname{Ric}|^2,\quad
 P=\int|\operatorname{Ric}^{\circ}|^2.
\]

When R\textgreater0, the project defect and its conformally useful
companion are

\[
 \mathcal E=4-\frac{S_1^2}{VR},\qquad
 J=\frac{4P}{R}=4-\frac{S_2}{R},\qquad
 \mathcal E-J=\frac{S_2-S_1^2/V}{R}.                 \tag{1.1}
\]

Assume every conformal class has negative Yamabe sign, as follows from
the actual positive-projection input used in I41. In particular
R\textgreater0 for every metric. Let alpha=inf\_g E\textless1. The
conformal reduction in I41, Section 3, gives inf\_g J=alpha.
Consequently the following \textbf{sharp positivity threshold} is an
exact reformulation of that infimum:

\[
 \mathcal Q_{\tau_*}(g):=R+\tau_*S_2\ge0\quad
 \hbox{for every smooth }g,\qquad
 \tau_*=-\frac1{4-\alpha}.                            \tag{1.2}
\]

For each tau\textless tau\_* there is a smooth metric for which
Q\_tau\textless0. Indeed, R/S\_2=1/(4-J) whenever S\_2\textgreater0, and
a sequence J-\textgreater alpha approaches the right-hand side in (1.2).
S\_2 cannot vanish identically in this negative Yamabe setting. No
minimizer is used in this argument.

Using the project Weyl-operator norm, Chern--Gauss--Bonnet writes the
same functional as

\[
 \mathcal Q_{\tau_*}
 =2\int|W|^2+
   \frac{1-\alpha}{3(4-\alpha)}S_2-16\pi^2\chi.       \tag{1.3}
\]

Thus an optimal defect is equivalent to an optimal positivity parameter
of an existing quadratic curvature functional. This does not imply that
every quadratic critical metric minimizes E, or that known rigidity
theorems apply in the parameter range (-1/3,-1/4).

\hypertarget{what-smooth-attainment-would-force}{%
\subsubsection{1.1. What smooth attainment would
force}\label{what-smooth-attainment-would-force}}

\textbf{Proposition 1.} In the preceding setting, a smooth metric
attaining alpha\textless1 has negative constant scalar curvature,
satisfies

\[
 \nabla\mathcal Q_{\tau_*}=0,\qquad
 \mathcal Q_{\tau_*}=0,                              \tag{1.4}
\]

and is a global minimizer of Q\_tau*. If 0\textless alpha\textless1,
this metric cannot be Bach-flat, Kähler for either orientation, or have
parallel Ricci tensor.

\textbf{Proof.} If E(g)=alpha, then
alpha\textless=J(g)\textless=E(g)=alpha by (1.1)--(1.2). The scalar
variance is therefore zero. Its constant value is negative because the
conformal class is negative. Equality in (1.2) gives Q\_tau*=0; global
nonnegativity implies stationarity and nonnegative second variation.

Here are direct verifications of the structural exclusions. They do not
require a general classification of quadratic critical metrics.

\begin{enumerate}
\def\labelenumi{\arabic{enumi}.}
\item
  If g is Bach-flat, the first variation of the Weyl energy vanishes. At
  constant scalar s, the gradient of S\_2 is -2s Ric\^{}circ. The
  coefficient of S\_2 in (1.3) is positive for alpha\textless1. Hence
  (1.4) gives Ric\^{}circ=0, which would give alpha=0. In particular
  this excludes either half-conformally-flat possibility.
\item
  For a Kähler surface in its complex orientation,
  R=S\_2/2-8pi\^{}2(2chi+3tau). Hence, along Kähler variations with the
  complex structure fixed,
  Q\_tau\emph{=(1/2+tau\_})S\_2-8pi\^{}2(2chi+3tau). The tensor
  h=Ric\^{}circ is an admissible such variation: the associated real
  (1,1)-form is the Ricci form minus (s/4) times the Kähler form, and is
  closed because s is constant. For sufficiently small real t, the
  perturbed form remains positive. Its derivative of S\_2 is -2s P,
  nonzero if P\textgreater0. Since 1/2+tau\_*\textgreater0, (1.4) forces
  P=0. Reversing the orientation changes the topological constant but
  not this calculation.
\item
  If Ric is parallel, put q=\textbar Ric\textbar\^{}2, constant on M.
  The classical first variation of R reduces to -2 Ric\^{}2+(q/2)g. This
  follows as well by varying the parallel factors: Rm(Ric)=Ric\^{}2.
  Also S\_2=s\^{}2V and Q\_tau\emph{=0 imply q+tau\_}s\^{}2=0. Equation
  (1.4) therefore reduces to Ric\^{}2+tau\_*s Ric=0. Each nonzero Ricci
  eigenvalue is q/s. If its rank is r, then q=s\^{}2/r and E=4-r. An
  E-value strictly between zero and one is impossible.
\end{enumerate}

These three contradictions prove the proposition. QED.

For reference, the ambient quadratic first-variation formulas and the
constant-scalar trace equation are classical; see
\href{https://www.mate.polimi.it/wp-content/themes/mymate/send_new.php?file=qdd205\&tipo=1}{Catino,
Section 2}. The passage from the project infimum to the zero level in
(1.4) uses I41. Catino's nonpositive-sectional-curvature theorem assumes
tau\textgreater=-1/4; it is not a theorem covering the
tau\_*\textless-1/4 case at issue here.

\hypertarget{a-stationary-version-and-normalization-check}{%
\subsubsection{1.2. A stationary version and normalization
check}\label{a-stationary-version-and-normalization-check}}

For completeness, the same Euler--Lagrange connection holds for any
smooth E-critical metric with S\_1\textless0, without assuming global
minimality. Under g\_t=e\^{}\{2tf\}g and Delta=d*d,

\[
 \dot V=4\int f,\quad \dot S_1=2\int sf,\quad
 \dot S_2=12\int s\Delta f,\quad
 \dot R=4\int s\Delta f.                             \tag{1.5}
\]

The last identity follows from the four-dimensional
Weyl/Chern--Gauss--Bonnet formula, not an assumption of conformal
invariance of Ricci energy. Stationarity of S\_1\^{}2/(VR) gives

\[
 \Delta(s-\bar s)-\frac{R}{S_1}(s-\bar s)=0.
\]

Pairing with s-bar s shows s is constant, since R/S\_1\textless0. At
this metric, the full first variation gives

\[
 \nabla R+\frac{2R}{sV}\operatorname{Ric}^{\circ}=0.
\]

It is precisely the critical equation for Q\_tau with
tau=-R/(s\^{}2V)=-1/(4-E), and Q\_tau=0. Conversely, a negative-csc
critical point of Q\_tau at its zero level, with tau\textless0, is
E-critical. This equivalence is a variational identity, not a new
existence theorem.

\hypertarget{consequences-for-the-actual-surgery-families}{%
\subsubsection{1.3. Consequences for the actual surgery
families}\label{consequences-for-the-actual-surgery-families}}

For the complex-hyperbolic handle-only family of I44,
0\textless h\textless1/3, alpha=4h/(1+h) and tau\_*=-(1+h)/4. Its sharp
quadratic inequality is

\[
 R-\frac{1+h}{4}S_2\ge0.                            \tag{1.6}
\]

This is also checked without the variational reduction: the exact I44
signature identity gives

\[
 R-\frac{1+h}{4}S_2
 =\frac{1-3h}{12}(S_2-32\pi^2A)+4\int|W^-|^2\ge0.
\]

I44 already proves nonattainment; (1.6) is not counted as a second
sharpness theorem. On the balanced family, only c
h\textless=alpha\textless=4h/(1+h) is known. Proposition 1 says that any
smooth attained subunit minimum there would have to be genuinely
non-Kähler, non-Bach-flat and non-parallel-Ricci. It neither determines
alpha nor rules out such a minimizer.

The remaining coefficient-four target has a useful invariant
cancellation. With D=(3/2)S\_2+18W\_+-72pi\^{}2p and T=A(1-h),

\[
 \frac89D+32\pi^2(p-A)-\Sigma
 =\frac13S_2+16W_+-32\pi^2A
 =4P-32\pi^2Ah.                                    \tag{1.7}
\]

The selected projection p cancels exactly. Thus merely selecting a
larger p, or averaging the same scalar projection inequality over more
sign classes, cannot by itself close this target. Any use of joint
actual classes must supply additional differential information.

\hypertarget{directional-torsion-rather-than-the-size-of-a-selected-projection}{%
\subsection{2. Directional torsion, rather than the size of a selected
projection}\label{directional-torsion-rather-than-the-size-of-a-selected-projection}}

The cancellation (1.7) identifies a geometric route that is different
from optimizing the old scalar-slack constant. Here the extra input is a
genuine global compatible almost complex structure, not a normalized
spinor defined only off its zeros.

Let g have scalar curvature -4 and let omega be a smooth self-dual
fundamental form with \textbar omega\textbar\^{}2=2. Define its Lee form
by d omega=theta wedge omega and put

\[
 \Theta=\int|\theta|^2,\qquad
 \nu=\int\bigl(|\nabla\omega|^2-|\theta|^2\bigr).
\]

The integrand in nu is nonnegative in dimension four. To see this
directly, complete omega locally to quaternionic self-dual forms
omega,omega\_2,omega\_3, each of norm squared two, and write nabla
omega=a tensor omega\_2+b tensor omega\_3. Put alpha=\emph{(a wedge
omega\_2), beta=}(b wedge omega\_3). The two maps defining alpha and
beta are isometries on one-forms. Therefore

\[
 |\nabla\omega|^2=2(|\alpha|^2+|\beta|^2),\qquad
 |\theta|^2=|d\omega|^2=|\alpha+\beta|^2,
\]

and the difference is \textbar alpha-beta\textbar\^{}2. Its vanishing is
the Hermitian intrinsic-torsion condition; equivalently the Nijenhuis
tensor is zero. This last assertion also follows by inserting the same
frame into the definition of N\_J. No particular normalization of
\textbar N\_J\textbar{} is needed below.

\textbf{Proposition 2 (directional energy bound).} Suppose the actual
scalar budget is S\_2\textgreater=32pi\^{}2A with A\textgreater0, and
T=2chi+3tau=A(1-h). Then

\[
 \boxed{
 P\ge4V-8\pi^2T-2\nu+2\Theta
                   +\frac{3(\nu-\Theta)^2}{8V}
 \ge8\pi^2Ah-2\nu+2\Theta
                   +\frac{3(\nu-\Theta)^2}{8V}.}       \tag{2.1}
\]

Only the second inequality uses the scalar budget. In particular,
nu=o(Ah) on a normalized low-defect sequence gives the missing
coefficient four lower estimate. More generally, if h\_j-\textgreater0,
R\_j/(8pi\^{}2A\_j)-\textgreater1, and
E(g\_j)/h\_j-\textgreater c\textless4, then every compatible global
omega\_j satisfies

\[
 \liminf_j\frac{\nu_j-\Theta_j}{A_jh_j}
                  \ge\pi^2(4-c).                    \tag{2.2}
\]

\textbf{Proof.} Set Q=\textless W\^{}+ omega,omega\textgreater. The
integrated self-dual Weitzenboeck identity and \textbar d
omega\textbar\^{}2=\textbar theta\textbar\^{}2 give

\[
 \int Q=-\frac43V+\frac12(\nu-\Theta).
\]

For a trace-free symmetric endomorphism of a three-dimensional Euclidean
space, \textbar W\^{}+\textbar\^{}2\textgreater=(3/8)Q\^{}2 when
\textbar omega\textbar\^{}2=2. Cauchy--Schwarz therefore yields

\[
 \int|W^+|^2\ge\frac{3}{8V}\left(\int Q\right)^2
 =\frac23V-\frac12(\nu-\Theta)
                  +\frac{3(\nu-\Theta)^2}{32V}.
\]

The chiral Chern--Weil identity is
P=S\_2/12+4int\textbar W\^{}+\textbar\^{}2-8pi\^{}2T. Since S\_2=16V,
substitution proves the first inequality in (2.1). The scalar budget
says V\textgreater=2pi\^{}2A, proving the second. Finally E=J=4P/R on
these constant-scalar metrics, so the assumed limits give
P/(Ah)-\textgreater2pi\^{}2c. Drop the nonnegative square in (2.1) and
rearrange to obtain (2.2). QED.

The Hermitian torsion formula is classical; a primary formula with the
same Lee convention is given in
\href{https://dergipark.org.tr/tr/download/article-file/2984722}{On Four
Dimensional Hermitian Manifolds, Proposition 3.1}. The inequalities
above are the displayed direct computation from that four-dimensional
decomposition and the scalar/chiral budgets.

\hypertarget{an-all-hermitian-metric-bound}{%
\subsubsection{2.1. An all-Hermitian-metric
bound}\label{an-all-hermitian-metric-bound}}

\textbf{Corollary 3.} Suppose that for every conformal class in question
there is the actual scalar input with the same A\textgreater0, and
T=A(1-h), 0\textless=h\textless=1/3. Every Hermitian metric in the
specified orientation satisfies

\[
 \boxed{\mathcal E(g)\ge J(g)\ge\frac{4h}{1+h}.}     \tag{2.3}
\]

\textbf{Proof.} The negative Yamabe representative, normalized to scalar
-4, is still Hermitian for the same complex structure. Thus nu=0.
Equation (2.1) implies P\textgreater=4V-8pi\^{}2T. The function

\[
 \frac{4(4V-8\pi^2T)}{8V-8\pi^2T}
\]

is increasing in V when T\textgreater0, and V\textgreater=2pi\^{}2A
makes its denominator positive and its numerator nonnegative. As
J=4P/(P+4V) is increasing in P, this proves J\_hat\textgreater=4h/(1+h).
The exact conformal formula in I41 says J(g) is a weighted average of
J\_hat and 1. The bound in (2.3) is at most one, so it transfers to g.
Equation (1.1) completes the proof. QED.

This covers a natural geometric class and uses no Lee-form smallness. It
does not assert that a balanced surgery manifold is Hermitian, or even
almost complex. If a compatible form is supplied only away from spinor
zeros, (2.1) cannot be applied by deleting those zeros and omitting the
boundary terms. Producing a global form with nu=o(Ah), or controlling
the omitted marked-core terms, remains a substantive extra requirement.

For comparison with the classical parameter inequalities, direct
expansion of D\_8 gives

\[
 \frac13S_2+16W_+=\frac8{27}D_8+\frac{64}{3}\pi^2p.  \tag{2.4}
\]

D\_8\textgreater=0 alone reaches only two thirds of the desired
32pi\^{}2A level when p=A. At the complex-hyperbolic calibration
D\_8=36pi\^{}2A, not zero. Accordingly the known non-saturation of D\_8
does not automatically determine the required first-order variation. The
new information in (2.1) is directional differential control, not an
increased projection.

\hypertarget{the-same-criterion-tolerates-isolated-projective-singularities}{%
\subsubsection{2.2. The same criterion tolerates isolated projective
singularities}\label{the-same-criterion-tolerates-isolated-projective-singularities}}

The smooth global almost complex hypothesis in Proposition 2 is stronger
than necessary. Its proof holds for any

\[
 \omega\in W^{1,2}(\Lambda^+),\qquad |\omega|^2=2
 \quad\hbox{almost everywhere}.                       \tag{2.5}
\]

Indeed multiplication by omega identifies one-forms with three-forms
isometrically almost everywhere, so d omega=theta wedge omega defines an
L\^{}2 Lee form. The quaternionic-frame calculation is pointwise algebra
on the weak derivative and gives nu\textgreater=0. The integrated
Weitzenboeck identity extends by smooth approximation in W\^{}\{1,2\};
the approximants need not have constant length. All curvature
coefficients are smooth on the fixed closed manifold, so every term
passes to the limit. Thus (2.1) and (2.2) remain valid with no
removed-boundary term.

For example, let a smooth positive spinor have only finitely many
transverse zeros. Away from these zeros its normalized projective
self-dual form has norm squared two. At a transverse zero,
\textbar Phi\textbar{} is comparable to distance r and \textbar nabla
Phi\textbar{} is bounded. Therefore its normalized form has
\textbar nabla omega\textbar\textless=C/r. This is square integrable in
four dimensions. The form extends weakly across the point: an annular
cutoff has squared gradient integral O(r\^{}2), tending to zero, and the
form is bounded. Consequently (2.5) holds even if the smooth manifold
has no smooth global almost complex structure.

This removes the isolated-transverse-zero obstruction to \textbf{using
the criterion}, not to proving its smallness premise. Actual spinor zero
sets have not been proved to be transverse here. Codimension-two
projective singularities need not have finite W\^{}\{1,2\} energy. Nor
is nu=o(Ah) a consequence of transversality alone. A weak unit form is
not being promoted to a smooth complex structure.

\hypertarget{a-completely-smooth-test-with-no-normalization-at-its-zeros}{%
\subsubsection{2.3. A completely smooth test with no normalization at
its
zeros}\label{a-completely-smooth-test-with-no-normalization-at-its-zeros}}

There is a second version that accepts every nonzero smooth self-dual
two-form omega. Define

\[
 a=\frac12\int|\omega|^2,\quad b=\frac14\int|\omega|^4,
 \quad V_{\rm eff}=\frac{a^2}{b}\le V,\quad
 T_\omega=\int\bigl(|\nabla\omega|^2-2|d\omega|^2\bigr).
\]

On a scalar -4 metric with the same actual budget,

\[
 \boxed{
 P\ge8\pi^2Ah+4(V-2\pi^2A)
   -\frac83(V-V_{\rm eff})-2\frac ab T_\omega
   +\frac{3T_\omega^2}{8b}.}                         \tag{2.6}
\]

To prove this, the Weitzenboeck identity gives
int\textless W\^{}+omega,omega\textgreater=T\_omega/2-4a/3. Also
\textbar(omega tensor
omega)\_0\textbar\^{}2=(2/3)\textbar omega\textbar\^{}4, so
Cauchy--Schwarz yields

\[
 \int|W^+|^2\ge\frac3{8b}(T_\omega/2-4a/3)^2
 =\frac23V_{\rm eff}-\frac{a}{2b}T_\omega
                       +\frac{3T_\omega^2}{32b}.
\]

Insert this into chiral Chern--Weil to obtain (2.6). No derivative of a
normalized form occurs. Thus V-V\_eff=o(Ah), (T\_omega)\_+=o(Ah), and
bounded a/b are sufficient for coefficient four on the normalized
low-defect sequences, including forms with arbitrary zeros.

For the ordinary actual unweighted spinor on scalar -4, take
omega=iF\_A\^{}+ and rho=\textbar Phi\textbar\^{}2\textless=4. Then
\textbar omega\textbar\^{}2=rho\^{}2/8, so b\textless=a and
V\_eff\textgreater=a. If K=int\textbar nabla\_A Phi\textbar\^{}2 and
Q=int(4-rho)\^{}2, the ordinary Weitzenboeck identity gives

\[
 V-V_{\rm eff}\le V-a
       =\frac1{16}\int(16-\rho^2)=\frac K2+\frac Q{16}. \tag{2.7}
\]

The existing actual scalar-slack estimates bound these terms at order
Ah, not at o(Ah); they also do not yet give the needed directional
estimate for T\_omega. Hence (2.6) explains a precise remaining loss
without claiming that the supplied smooth form already closes it.

\hypertarget{finite-defect-estimate-with-a-directional-cost}{%
\subsubsection{2.4. Finite defect estimate with a directional
cost}\label{finite-defect-estimate-with-a-directional-cost}}

If the actual beta=2 budget is also available, then
R\textgreater=8pi\^{}2A(1+h), as in I41. Set t\_omega=nu-Theta for a
unit form in (2.5). For 0\textless=h\textless=1/3, (2.1) and
V\textgreater=2pi\^{}2A give

\[
 P-4hV\ge4(1-h)(V-2\pi^2A)-2t_\omega
                                    +\frac{3t_\omega^2}{8V}.
\]

Consequently, on the normalized scalar -4 representative,

\[
 \boxed{\mathcal E\ge\frac{4h}{1+h}
 -\frac{[t_\omega-3t_\omega^2/(16V)]_+}
             {\pi^2A(1+h)^2}
 \ge\frac{4h}{1+h}
 -\frac{(\nu-\Theta)_+}{\pi^2A(1+h)^2}.}             \tag{2.8}
\]

Indeed E-4h/(1+h)=4(P-4hV)/{[}(1+h)R{]}, and the displayed R lower bound
controls the negative part. Thus the sufficient directional condition is
actually (nu-Theta)\_+=o(Ah), weaker than nu=o(Ah). Formula (2.8) is a
finite estimate with an independent first-derivative cost; it is not an
estimate asserting that this cost vanishes for the actual projective
form.

For a nonzero smooth form as in Section 2.3, the scale-invariant
replacement for this cost is

\[
 C_\omega=\frac43(V-V_{\rm eff})+\frac abT_\omega
                         -\frac{3T_\omega^2}{16b}.
\]

The identical argument gives

\[
 \boxed{\mathcal E\ge\frac{4h}{1+h}
       -\frac{[C_\omega]_+}{\pi^2A(1+h)^2}.}         \tag{2.9}
\]

Indeed (2.6) gives P-4hV\textgreater=4(1-h)(V-2pi\^{}2A)-2C\_omega.
Multiplication of omega by a nonzero constant leaves C\_omega unchanged.
For the ordinary actual form, harmonic projection gives
a\textgreater=2pi\^{}2p, while b\textless=a and Cauchy--Schwarz give

\[
 1\le a/b\le V/a\le V/(2\pi^2A),\qquad
 V-V_{\rm eff}\le V-2\pi^2p.
\]

Thus a/b tends to one on the normalized low-defect sequences; its
boundedness is not a further premise there. The smallness of the signed
directional cost is still unproved. In particular t\_omega itself can be
rewritten by Weitzenboeck as a directional Weyl contraction; the formula
does not secretly supply new differential control of that contraction.

\hypertarget{the-bounded-ricci-euler-density-endpoint-gap}{%
\subsection{3. The bounded-Ricci Euler-density endpoint
gap}\label{the-bounded-ricci-euler-density-endpoint-gap}}

The previously open bounded-Ricci version of the Euler-density bridge
can be closed using an explicit classical extension that was not used in
I43. The unrestricted weak-Ricci problem is different and remains open.

For a bounded self-dual form define the nonnegative density

\[
 Z_t=|\operatorname{Ric}+g|^2+
 \left|W^+-t\left(\frac{|\omega|^2}{3}I-\omega\otimes\omega\right)\right|^2
 +(|\omega|^2-2)^2+|d\omega|^2,
 \qquad 0\le t\le\tfrac14.                          \tag{3.1}
\]

\textbf{Theorem 4.} Let M be a closed oriented smooth four-manifold,
omega a smooth self-dual form, and t a constant in {[}0,1/4{]}. Given
finite bounds K,L,Lambda\_0, there is c\textgreater0 depending only on
these bounds such that

\[
 |\operatorname{Ric}|\le K,\quad |\omega|\le L,\quad
 \frac{8\pi^2\chi(M)}{V}\le\Lambda_0
 \quad\Longrightarrow\quad
 \boxed{\fint_M Z_t\ge c(K,L,\Lambda_0)>0.}           \tag{3.2}
\]

No Einstein equation, diameter, volume lower bound, global Sobolev
bound, or all-center small-energy hypothesis is assumed. The Ricci bound
is pointwise and is \textbf{not} removed by this theorem.

\hypertarget{the-precise-classical-input}{%
\subsubsection{3.1. The precise classical
input}\label{the-precise-classical-input}}

\href{https://doi.org/10.1090/S0894-0347-05-00511-4}{Cheeger--Tian,
Theorem 0.8 and Remark 8.22} give the following bounded-Ricci local
input. If \textbar Ric\textbar\textless=K and the unnormalized curvature
L\^{}2 energy on B\_r(x), r\textless=1, is less than a threshold
epsilon\_CT(K), an interior neighborhood has uniformly bounded
C\^{}\{1,alpha\} and W\^{}\{2,q\} covering geometry. The bounded-Ricci
conclusion is not a pointwise full-curvature estimate. Fix one finite
q\textgreater4.

The source was checked in full at the relevant printed pages: Theorem
0.8 on p.489, Remark 8.22 on p.515, and the meaning of covering geometry
in Remarks 2.7 and 4.8 on pp.501 and 507. The original article is also
available in this
\href{https://www.math.stonybrook.edu/Courses/MAT683/201001/MAT683S10.pdf}{institutional
course collection}. Remark 4.8 is used to interpret the local coverings,
not as an unconditional fixed-scale regularity theorem. The
bounded-Ricci use of Theorem 0.8 together with Remark 8.22 is also made
in \href{https://math.berkeley.edu/~lott/4d3.pdf}{Lott, Section 3.1}.

In particular, rescaling by r\^{}\{-2\} supplies constants
a,Q\textgreater0, depending only on K, and a local diffeomorphism

\[
 F_x:B_{\mathbb R^4}(0,ar)\longrightarrow M,\qquad F_x(0)=x,
 \qquad Q^{-1}g_{\rm Euc}\le F_x^*g\le Qg_{\rm Euc}. \tag{3.3}
\]

This is a harmonic coordinate chart upstairs followed by the genuine
covering projection. It need not be injective downstairs. It is a local
isometry for gamma\_x=F\_x\^{}*g. The fields and their errors remain
those of the original g; no smoothing of the actual equations is
asserted. The chart retains the fixed C\^{}\{1,alpha\} and W\^{}\{2,q\}
bounds, not only the ellipticity displayed in (3.3).

\hypertarget{local-isometries-suffice-for-the-averaging-argument}{%
\subsubsection{3.2. Local isometries suffice for the averaging
argument}\label{local-isometries-suffice-for-the-averaging-argument}}

We use the following extension of I43. Suppose that every x in a
measurable set G has a pointed local isometry
F\_x:(N\_x,gamma\_x,z\_x)-\textgreater(M,g,x), with the 4rho ball
compactly contained in N\_x and Vol\_gamma B\_rho(z\_x)\textgreater=v
rho\^{}4, for fixed rho,v\textgreater0. Then

\[
 \frac{\operatorname{Vol}(G)}V
               \le C(K,L,\rho,v)\fint_M Z_t.         \tag{3.4}
\]

Indeed the buffered local coercivity proved in I42--I43 applies upstairs
and gives
int\_\{B\_\{2rho\}(z\_x)\}F\_x\^{}*Z\_t\textgreater=epsilon(K,L,rho,v)\textgreater0.
In polar coordinates, integrate only up to the gamma\_x cut time. The
Ricci bound controls the radial Jacobian by C\_J(K,rho)s\^{}3. Drop the
cut-time restriction after making this estimate. Projected geodesics are
original-g geodesics, so

\[
 \epsilon\le C_J\int_{S_xM}\int_0^{2\rho}
           s^3 Z_t(\exp_x^g(s\theta))\,ds\,d\theta.  \tag{3.5}
\]

This pointwise lower bound is now expressed without the choice of local
chart. Integrate over G, enlarge to M by nonnegativity, and use
invariance of original-g Liouville measure to obtain

\[
 \epsilon\operatorname{Vol}(G)
 \le C_J|S^3|\frac{(2\rho)^4}{4}\int_M Z_t.
\]

This proves (3.4). Neither the number of sheets nor measurable selection
of a family of charts enters. An exponential map is not declared to be a
covering map by this argument.

For (3.3), take rho=ar/(8sqrt Q). Uniform ellipticity gives
B\_\{4rho\}\^{}\{gamma\}(0) contained in the Euclidean ball of radius
ar/2, and

\[
 \operatorname{Vol}_{\gamma}B_\rho(0)
                         \ge\omega_4Q^{-4}\rho^4.
\]

Thus the small-energy charts satisfy precisely the buffered hypotheses.

For these charts the local coercivity also has a direct proof, avoiding
any completeness convention for the upstairs domain. If their local Z\_t
integrals tend to zero, harmonic coordinates give C\^{}\{1,alpha'\} and
weak W\^{}\{2,q\} subsequential convergence on smaller charts. Bounded
Ricci and small L\^{}2 Ric+g imply small L\^{}q Ric+g for every fixed
finite q, by interpolation. The harmonic-coordinate metric equation
makes the limit Einstein and smooth and upgrades convergence to strong
interior W\^{}\{2,q\}. Interior d+delta estimates give an H\^{}1 bound
for the self-dual forms; Rellich and their L-infinity bound give strong
convergence of their quadratic models. The limit has Ric=-g, a closed
self-dual form of squared length two, and W\^{}+=t(2I/3-omega tensor
omega). At t=0 Armstrong's local ASD Einstein almost-Kaehler obstruction
applies. At t\textgreater0 the Einstein Bianchi identity forces the form
parallel, hence t=1/2, outside the interval. This is the same local
contradiction as I42, now directly in the already regular covering
charts.

\hypertarget{selecting-a-positive-volume-fraction-of-good-centers}{%
\subsubsection{3.3. Selecting a positive volume fraction of good
centers}\label{selecting-a-positive-volume-fraction-of-good-centers}}

In the curvature-operator convention,

\[
 \int|\mathscr R|^2=8\pi^2\chi+P,
 \qquad \fint|\mathscr R|^2\le
 E_*:=1+K^2+\max\{\Lambda_0,0\}.                    \tag{3.6}
\]

For the four-index Riemann tensor the identity is instead
int\textbar Rm\textbar\^{}2=32pi\^{}2chi+4P; changing conventions only
changes the local energy threshold by four. Bishop comparison gives Vol
B\_r(y)\textless=C\_B(K)r\^{}4 for r\textless=1. Therefore Fubini gives

\[
 \fint_M\left(\int_{B_r(x)}|\mathscr R|^2\right)d\mu(x)
                          \le C_B(K)r^4E_*.
\]

Choose

\[
 r=\min\left\{1,
    \left[\frac{\epsilon_{\rm CT}(K)}{2C_B(K)E_*}\right]^{1/4}\right\}.
\]

Markov's inequality shows that G=\{x:int\_\{B\_r(x)\}\textbar mathscr
R\textbar\^{}2\textless=epsilon\_CT(K)\} has volume at least V/2. At
every x in G, (3.3) supplies a chart with the same a,Q,rho,v. Equation
(3.4) now proves (3.2). QED.

In contrast to the earlier global Li-II smoothing argument, this proof
does not require every center to be low energy. The exceptional centers
may retain curvature concentration. No pointwise curvature bound is
deduced there or on general bounded-Ricci metrics.

\hypertarget{actual-gauge-and-topology-only-ratio-inputs}{%
\subsubsection{3.4. Actual gauge and topology-only ratio
inputs}\label{actual-gauge-and-topology-only-ratio-inputs}}

Use the same bounded actual form supplied in I06/I09, with t=1/beta,
beta\textgreater=4. Its estimate is avg
Z\_t\textless=C(r\_Ric+D\_beta/V)+o(1), where
r\_Ric=avg\textbar Ric+g\textbar\^{}2. Sending the freely selectable
auxiliary allowance to zero after applying (3.2) gives

\[
 \boxed{r_{\rm Ric}+D_\beta/V\ge c(K,\Lambda_0),
                           \qquad\beta\ge4.}         \tag{3.7}
\]

The actual structure remains fixed through all weights used to supply
that form. No square-cost Hessian theorem is inserted into this
zero/first-order form argument.

An absolute volume premise is unnecessary if
p\textgreater=A\textgreater0 and chi/A\textless=B. For
r\_Ric\textless=1, the actual scalar estimate and the trace inequality
imply

\[
 32\pi^2A\le S_2\le (4+2\sqrt{r_{\rm Ric}})^2V,
 \qquad \frac{8\pi^2\chi}{V}\le9\max\{B,0\}.        \tag{3.8}
\]

For r\_Ric\textgreater1, nonnegativity of D\_beta already gives a fixed
gap. Hence (3.7) holds with c(K,B) under this topological ratio
hypothesis. If
delta\_-=\textbar\textbar Ric\^{}sharp+I\textbar\textbar{}\emph{\{infinity,op\}\textless=1/4,
the pointwise Ricci bound is uniform and r\_Ric\textless=4delta}-\^{}2.
Thus

\[
 \boxed{\delta_-^2+D_\beta/V\ge c(B)>0,
                              \qquad\beta\ge4.}      \tag{3.9}
\]

In the existing upper-wall setting, retain p\textgreater=A\textgreater0,
T\_+=A(1-h), h in {[}1/4,1/3{]}, e=R/(8pi\^{}2A), and all actual input
quantifiers of I43. Its unchanged algebra, with
G(h)=(4+sqrt(1-3h))\^{}2/9, now gives

\[
 \boxed{\delta_-^2+(G(h)-e)_+\ge c(B)>0.}             \tag{3.10}
\]

The bounded-diameter premise and the special elliptic-metric premise of
the earlier versions are no longer needed. This is an upper-energy
exclusion, not a claim that topology alone places every candidate near
that upper wall. In particular it is not a full Yau26 existence theorem.

For the actual complex-hyperbolic surgery family,
chi/A=1/3+k/A-2ell/A\textless=1/3+h\textless=2/3. Thus its
pointwise-input conclusions above have a uniform constant without any
additional Euler-ratio assumption. This is a direct application, not a
new gauge supplier.

Finally, any sequence with bounded Ricci and bounded form amplitude,
t\_j in {[}0,1/4{]}, and avg Z\_\{t\_j\}-\textgreater0 must satisfy

\[
 \frac{8\pi^2\chi(M_j)}{V_j}\longrightarrow+\infty.   \tag{3.11}
\]

Otherwise a subsequence has a finite upper density bound and contradicts
(3.2). The same conclusion holds for the actual input of (3.7).
Constants remain non-effective because the inherited local coercivity is
proved by compactness. No uniform pointwise Ricci bound has been derived
from the original weak integral defect.

\hypertarget{the-entire-finite-pencil-with-improved-zero-end-decay}{%
\subsubsection{3.5. The entire finite pencil, with improved zero-end
decay}\label{the-entire-finite-pencil-with-improved-zero-end-decay}}

Retain the exact conventions of I07 and I09:

\[
 b_a=\frac{(-a^2)_+}{V},\qquad
 D_\beta=\frac{\beta+1}{\beta}
              \int(s^2+6\beta|W^+|^2)-72\pi^2p,
 \qquad \kappa_\beta=r_{\rm Ric}+D_\beta/V.
\]

In particular the definition of b\_a contains no 8pi\^{}2 factor. Under
the actual solvability of I07 and the geometric hypotheses of (3.2),

\[
 \boxed{(1+b_a)\kappa_\beta\ge
 c(K,\Lambda_0)\frac{(\beta-2)^2}{\beta^2},
                      \qquad 2\le\beta<\infty.}     \tag{3.12}
\]

For 2\textless=beta\textless=4 the inherited finite-pencil theorem with
cutoff L=4 gives (1+b\_a)kappa\_beta\textgreater=c\_4(beta-2)\^{}2,
hence at least 4c\_4(beta-2)\^{}2/beta\^{}2. For beta\textgreater=4, use
(3.7) and (beta-2)\^{}2/beta\^{}2\textless=1. The smaller of these two
constants proves (3.12). If a\^{}2\textgreater=T\_+, the inherited
compact-interval estimate absorbs b\_a; the tail already has no square
cost. Thus the factor 1+b\_a can be omitted under that hypothesis. As in
(3.8), p\textgreater=A and chi/A\textless=B replace the Euler-density
premise, with constant c(K,B).

For the compactified pencil P\_t=tD\_\{1/t\}, including
P\_0=6int\textbar W\^{}+\textbar\^{}2, the elementary inequality
r+tD/V\textgreater=t(r+D/V) gives

\[
 \boxed{(1+b_a)(r_{\rm Ric}+P_t/V)
       \ge c\,t(1-2t)^2,\qquad0\le t\le\tfrac12.}   \tag{3.13}
\]

The previous unrestricted estimate has t\^{}3 instead of t. This
improvement is proved only with the stated pointwise Ricci and
Euler-density control. At t=0 (3.13) itself is trivial; (3.2) is the
nontrivial form statement at that endpoint. No claim of an optimal
exponent is made.

\hypertarget{the-topology-itself-can-supply-the-missing-density-budget}{%
\subsubsection{3.6. The topology itself can supply the missing density
budget}\label{the-topology-itself-can-supply-the-missing-density-budget}}

The Euler-density condition can be replaced by a topological ratio,
without a preassigned monopole projection. Put T\_+=2chi+3tau and z=avg
Z\_t. Directly from (3.1) and chiral Chern--Weil,

\[
 \frac{4\pi^2T_+}{V}
  =\frac23+\frac{16t^2}{3}+O_L(\sqrt z+z).          \tag{3.14}
\]

To verify the error, the model tensor has squared norm
(2/3)\textbar omega\textbar\^{}4. Its average differs from 8/3 by
O\_L(sqrt z), since \textbar\textbar omega\textbar\^{}2-2\textbar{} has
small L2 norm and omega is bounded. The L2 model-fit error changes
avg\textbar W\^{}+\textbar\^{}2 by O\_L(sqrt z+z). Also avg
s\^{}2=16+O(sqrt z+z) and
avg\textbar Ric\^{}circ\textbar\^{}2\textless=z. Insert these into
4pi\^{}2T\_+=2int\textbar W\^{}+\textbar\^{}2+S\_2/24-P/2.

For z\textless=epsilon(L) the right side lies in {[}1/3,4/3{]}. In
particular T\_+\textgreater0 and
3pi\^{}2T\_+\textless=V\textless=12pi\^{}2T\_+. Thus, if
T\_+\textgreater0 and chi/T\_+\textless=B,

\[
 \frac{8\pi^2\chi}{V}\le\frac83\max\{B,0\}.
\]

Theorem 4 now gives z\textgreater=c(K,L,B)\textgreater0, after including
the alternative z\textgreater epsilon(L). If T\_+\textless=0, a stronger
direct bound is available: Chern--Weil gives P\textgreater=S\_2/12,
whence

\[
 r_{\rm Ric}=P/V+\fint(s+4)^2/4
 \ge\fint[s^2/12+(s+4)^2/4]
 =1+\tfrac13\fint(s+3)^2\ge1.                      \tag{3.15}
\]

In particular every fixed oriented smooth four-manifold has a
bounded-Ricci, bounded-form endpoint gap, without a diameter condition.
Its constant may depend on its topology. All actual consequences (3.7),
(3.12), and (3.13) hold with the density assumption replaced by
T\_+\textgreater0 and chi/T\_+\textless=B. This replacement is derived
in the small-defect regime; arbitrary metrics have not been assigned
topologically fixed volume.

\hypertarget{uniformity-on-the-full-classical-general-type-surgery-family}{%
\subsubsection{3.7. Uniformity on the full classical general-type
surgery
family}\label{uniformity-on-the-full-classical-general-type-surgery-family}}

Let X be any minimal complex surface of general type,
A=c\_1\^{}2(X)\textgreater0, and M=X\#k CPbar\^{}2\#ell(S\^{}1 times
S\^{}3), with h=(k+4ell)/A in {[}0,1/3{]}. Keep precisely the actual
polarized/chamber supply in I03, including its b\_+=1 and loop-insertion
quantifiers. The selected class has a\^{}2=A-k\textgreater=T\_+=A(1-h),
so b\_a=0.

The classical Noether inequality A\textgreater=2p\_g-4 and the
identities chi\_h=1-q+p\_g, chi(X)=12chi\_h-A imply

\[
 \chi(X)\le5A+36,\qquad \chi(X)/A\le41.
\]

Here A is a positive integer and q\textgreater=0. These formulas and the
full scope of Noether's inequality are recorded in
\href{https://arxiv.org/pdf/math/0702548}{Liedtke, Theorem 2.1} and
\href{https://www.mathe8.uni-bayreuth.de/de/team/prof-fabrizio-catanese-old/pdf/101.pdf}{Bauer--Catanese--Pignatelli,
p.5}. No regularity or ampleness premise is added to X. Surgery gives

\[
 \frac{\chi(M)}A\le41+h<42,\qquad
 \frac{\chi(M)}{T_+}\le\frac{41+h}{1-h}\le62.        \tag{3.16}
\]

Thus the Euler ratio in the preceding theorems is automatic across this
entire classical family, not just complex-hyperbolic bases. For every
metric on it and every finite beta\textgreater=2 there is a universal
c\textgreater0 such that

\[
 \boxed{\delta_-(g)^2+D_\beta(g,a_g)/V
             \ge c\frac{(\beta-2)^2}{\beta^2}.}     \tag{3.17}
\]

Indeed delta\_-\textless=1/4 supplies a universal pointwise Ricci bound
and r\_Ric\textless=4delta\_-\^{}2, so use (3.12), (3.16) and b\_a=0. If
delta\_-\textgreater1/4, the left side is at least 1/16 and the
right-hand parameter factor is at most one. Reduce c accordingly. The
class can be selected separately for each original metric, but remains
fixed through its auxiliary weights.

The upper-wall consequence (3.10), for h in {[}1/4,1/3{]}, also has a
universal constant across this family. There is no diameter, volume,
base noncollapse or separately imposed Euler budget. Equation (3.17)
uses the pointwise defect delta\_-, not the unrestricted integral E; it
therefore does not assert the still-open integral zero-end theorem.

In fact (3.17) extends to all nonnegative surgery counts, with no upper
bound on h. For h\textgreater=1/3 the already established I03
obstruction gives E\textgreater=epsilon\_0. Since
E\textless=4delta\_-\^{}2, the left side of (3.17) is at least
epsilon\_0/4, while its parameter factor is at most one. Combining the
two h ranges proves the same universal pointwise pencil inequality on
the entire eligible general-type blowup/handle family. This last step
uses the old critical-band theorem; it is not an additional zero-end
regularity assertion beyond (3.16).

\hypertarget{actual-annular-matching-without-a-first-order-ale-expansion}{%
\subsection{4. Actual annular matching without a first-order ALE
expansion}\label{actual-annular-matching-without-a-first-order-ale-expansion}}

This is a separate consequence of the actual prolonged equations. It
does not use Section 3's pointwise Ricci hypothesis, and does not assert
that the original weak-error sequence has an ALE decomposition.

Use the actual fields of I39/I41: Phi in S\^{}+, determinant connection
A, rho=\textbar Phi\textbar\^{}2\textless=4, P\_spin=nabla\_A Phi,
zeta=\textless Phi,P\_spin\textgreater, and the smooth cubic torsion
t\_Phi in E\_-=S\^{}- tensor L, L=det S\^{}+. Its induced determinant
connection is 3A. For beta\textgreater=4 put

\[
 \rho_\beta=\frac23\sqrt{4(1+1/\beta)},\quad
 U_\beta=\frac43(2/\beta-1),\quad
 m_\beta=-U_\beta\rho_\beta^3/2>0,
\] \[
 S=m_\beta+2\rho^2|P_{\rm spin}|^2,\quad
 d=2\rho^2|P_{\rm spin}|^2-|t_\Phi|^2\ge0,\quad
 \mathcal J=\frac{(S-2m_\beta)^2}{S},\quad
 \mathcal A=(\rho-\rho_\beta)^2+\mathcal J+d,
 \quad \mathcal R=D^-_{3A}t_\Phi.                    \tag{4.1}
\]

The old Dirac algebra gives d\textless=4rho\textbar zeta\textbar\^{}2.
There is no division by rho or t\_Phi in (4.1). Let
eta=kappa\_beta+epsilon\_aux\textless=1. Under the unchanged actual
supply and a\^{}2\textgreater=T\_+, the earlier proofs give

\[
 \int\mathcal A\le C\sqrt\eta V,\qquad
 \int|\mathcal R|^{4/3}\le C\eta^{2/3}V,             \tag{4.2}
\] \[
 \int_B|F_A|^2\le C\operatorname{Vol}(B)
                      +C(\beta^{-1}+\sqrt\eta)V     \tag{4.3}
\]

for every measurable B. The first two statements are from Sections 2 and
8 of the \href{20260927_mixed_current_and_signed_cores.md}{mixed-current
proof}. For (4.3), Section 6 of the
\href{20260927_conformal_separation_and_zero_end.md}{conformal
continuation} supplies the anti-self-dual part. The positive part
follows from \textbar F\_A\^{}+\textbar\^{}2=f\^{}2rho\^{}2/8, bounded
rho, and the actual weight estimate
int(f-sqrt(4(1+1/beta)))\^{}2\textless=C eta V. Thus (4.3) concerns the
full determinant curvature, not just one chirality.

\hypertarget{which-flat-quotient-holonomies-admit-both-limiting-fields}{%
\subsubsection{4.1. Which flat quotient holonomies admit both limiting
fields?}\label{which-flat-quotient-holonomies-admit-both-limiting-fields}}

Let Gamma be a finite subgroup of SO(4) acting freely on S\^{}3, with
its specified orientation and a flat Spin-c lift. Represent Spin\^{}c(4)
as

\[
 \{(A_+,A_-)\in U(2)\times U(2):\det A_+=\det A_-\}.
\]

These representations act on S\^{}+ and S\^{}- respectively, not on
E\_-. Suppose both S\^{}+ and E\_-=S\^{}- tensor det S\^{}+ have a
nonzero invariant vector. Choose these vectors first in their unitary
bases, and write lambda=det A\_+=det A\_-. Then

\[
 A_+=\operatorname{diag}(1,\lambda),\qquad
 A_-=\operatorname{diag}(\lambda^{-1},\lambda^2),
 \qquad E_-:\operatorname{diag}(1,\lambda^3).        \tag{4.4}
\]

The real Clifford space can be written as Q(z\_1,z\_2)=((z\_1,-bar
z\_2),(z\_2,bar z\_1)). Direct multiplication gives

\[
 A_-Q(z_1,z_2)A_+^{-1}
                   =Q(\lambda^{-1}z_1,\lambda^2z_2). \tag{4.5}
\]

The character lambda is faithful, since lambda=1 acts trivially on the
tangent space. Its image is finite in U(1), so Gamma is cyclic. Its
action is free precisely when its order n is odd: for even n the element
lambda=-1 fixes the second axis. Conversely all these odd cyclic actions
admit (4.4). Thus the oriented quotient types allowing both invariants
are exactly

\[
                         \frac1n(1,-2),\qquad n\text{ odd}.        \tag{4.6}
\]

Here n=1 is the trivial group. For the usual notation 1/n(1,q), let zeta
generate the roots of unity and delta=zeta\^{}\{1+q\}. The canonical
lift is A\_+\^{}0=diag(1,delta), A\_-\^{}0=diag(zeta,zeta\^{}q). A
character twist chi gives S\^{}+=chi A\_+\^{}0 and E\_-=chi\^{}3 delta
A\_-\^{}0. A positive invariant forces chi=1 or delta\^{}\{-1\};
checking the other representation then gives

\[
 q+2=0\pmod n\quad\hbox{or}\quad2q+1=0\pmod n.      \tag{4.7}
\]

They are the two coordinate orders of (4.6). In these cases the
permitted determinant characters are delta and delta\^{}\{-1\}. In the
fixed marking 1/n(1,-2), this is exactly
c\_1(L)\textbar{}\emph{\{S\^{}3/Gamma\}=+1 or -1 in Z/n. For odd n the
determinant class specifies the Spin-c structure, since multiplication
by two is invertible in Z/n.~For n=3 the invariant space in E}- is
two-dimensional, but the two allowed determinant classes are still
distinct. Orientation reversal is not an extra permitted change of
marking.

Call a marked group/lift \textbf{nonresonant} when it does not admit
these two invariant vectors. In particular every nontrivial Gamma in
SU(2)\_- is nonresonant. This definition is an explicit finite
representation condition, not a new analytical assumption hidden in
terminology.

The Spin-c presentation is classical; see
\href{https://haydys.net/misc/IntroGaugeTheory_LectNotes.pdf}{Haydys,
Example 121(b)}. Equations (4.4)--(4.7) are the direct representation
calculation used here.

\hypertarget{quantitative-annular-obstruction}{%
\subsubsection{4.2. Quantitative annular
obstruction}\label{quantitative-annular-obstruction}}

Let Q\_Gamma=\{1/2\textless\textbar y\textbar\textless4\}/Gamma carry
its flat cone metric h\_C. Fix a finite family of nonresonant marked
groups/lifts. There are delta\_0,c\_0\textgreater0 such that, for an
orientation-preserving embedded annulus iota:Q\_Gamma-\textgreater M and
r\textgreater0 with

\[
 \|r^{-2}\iota^*g-h_C\|_{C^2}<\delta_0,
\]

the actual fields restricted to that annulus satisfy

\[
 \boxed{c_0\le r^{-4}\int_{\iota Q_\Gamma}\mathcal A
          +\int_{\iota Q_\Gamma}|F_A|^2
          +r^{-8/3}\int_{\iota Q_\Gamma}|\mathcal R|^{4/3}.}       \tag{4.8}
\]

Only the actual Dirac equation, bounded spinor amplitude, and
definitions (4.1) are needed locally. No first-order ALE expansion,
finite energy of such an expansion, or relation between r and a smaller
core scale is assumed. Constants are not asserted uniform in arbitrary
quotient order.

\textbf{Proof.} If no such constants exist, take annuli with rescaled
metrics h\_j-\textgreater h\_C in C\^{}2 and right side tending to zero.
Pass to a fixed group, Spin-c type, and limit beta in the
compactification {[}4,infinity{]}. Under homothetic spinor
identification transport the physical t\_Phi without changing its
amplitude. This field equals r\^{}\{-1\}t\_\{Phi,h\}, not the cubic
field recomputed at scale h. Hence

\[
 D^+_{h,A}\Phi=0,\qquad
 D^-_{h,3A}\widetilde t=r\mathcal R,
\]

and the three terms in (4.8) become the integrals of A,
\textbar F\_A\textbar\^{}2, and \textbar D\^{}-\textbackslash widetilde
t\textbar\^{}\{4/3\} on the fixed annulus. S in that A still uses the
physical derivative, equivalently
r\^{}\{-2\}\textbar nabla\_\{h,A\}Phi\textbar\^{}2.

The determinant line is torsion on Q\_Gamma and admits a flat
connection. There are finitely many flat characters, so choose a
subsequence with one flat reference A\_0. Since
H\^{}1(Q\_Gamma;R)=H\^{}2(Q\_Gamma;R)=0, abelian Neumann Coulomb gauge
gives
\textbar\textbar A-A\_0\textbar\textbar\_\{W\^{}\{1,2\}\}\textless=C\textbar\textbar F\_A\textbar\textbar{}\emph{2-\textgreater0,
hence convergence in L\^{}4. The same holds for the induced twist in
E}-. There is no premise that the original determinant curvature
vanishes.

The positive spinors are bounded; their Dirac equation and interior
elliptic estimates give local W\^{}\{1,2\} bounds. Their squared norms
tend to rho\_beta\textgreater0 in L\^{}1. For the opposite spinor, first
observe

\[
 \mathcal J\ge S-4m_\beta,\qquad
 \int|S-2m_\beta|
      \le\left(\int\mathcal J\int S\right)^{1/2}.
\]

Thus S-\textgreater2m\_beta in L\^{}1 and \textbar widetilde
t\textbar\^{}2=S-m\_beta-d-\textgreater m\_beta in L\^{}1. In particular
\textbar widetilde t\textbar\^{}2 is uniformly integrable, also for the
varying volume forms. The Dirac estimate and
\textbar\textbar(A-A\_0)widetilde
t\textbar\textbar\_\{4/3\}\textless=\textbar\textbar A-A\_0\textbar\textbar\_4\textbar\textbar widetilde
t\textbar\textbar\_2 give local W\^{}\{1,4/3\} bounds. This embedding
into L\^{}2 is critical, not compact. Instead use subcritical Rellich to
obtain almost-everywhere convergence, then the preceding uniform
integrability and Vitali to obtain strong L\^{}2.

Both limits solve their flat Dirac equations and have constant positive
norm. Interior regularity makes them smooth. The flat Lichnerowicz
identity, applied locally to the constant norm, forces both to be
parallel. They therefore give the two holonomy invariants excluded by
the marked nonresonance condition. This contradiction proves (4.8). QED.

The same proof includes a resonant quotient with a forbidden determinant
class. For example, an annulus marked 1/n(1,-2) with odd n and
restricted c\_1(L) not equal to either +1 or -1 has the lower bound
(4.8).

\hypertarget{packing-bounds-at-the-actual-physical-scales}{%
\subsubsection{4.3. Packing bounds at the actual physical
scales}\label{packing-bounds-at-the-actual-physical-scales}}

Suppose the sequence has annuli as above, from a fixed finite marked
nonresonant family, with overlap at most Lambda and radii
r\_i\textless=r\_max. Take r\_max smaller than a fixed constant
depending on this data. Multiply (4.8) by r\_i\^{}4, sum, use
(4.2)--(4.3) on the union, and absorb C r\_max\^{}4 sum\_i r\_i\^{}4.
Since Vol(iota\_i Q\_i)\textless=C r\_i\^{}4, the result is

\[
 \boxed{\frac{\sum_i r_i^4}{V}\le C\Lambda
  [\sqrt\eta+r_{\max}^4(\beta^{-1}+\sqrt\eta)
                         +r_{\max}^{4/3}\eta^{2/3}].}            \tag{4.9}
\]

If each annulus additionally obeys r\_i\^{}\{-4\}int\_\{iota\_i
Q\_i\}A\textless=c\_0/2, absorb that term in (4.8), multiply by
r\_i\^{}\{8/3\}, and sum instead. One obtains

\[
 \boxed{\frac{\sum_i r_i^{8/3}}V\le C\Lambda
       [\eta^{2/3}+r_{\max}^{8/3}(\beta^{-1}+\sqrt\eta)].}       \tag{4.10}
\]

The determinant contribution has an absorbable part C
r\_max\^{}\{8/3\}sum\_i r\_i\^{}4\textless=C r\_max\^{}4 sum\_i
r\_i\^{}\{8/3\}. Thus no separate local determinant-smallness condition
was inserted.

For comparable radii r\_i comparable to r and N\textgreater=nu V, along
beta-\textgreater infinity and eta-\textgreater0 choose r\_max
comparable to that r. The terms with beta\^{}\{-1\}+sqrt eta can then be
absorbed. Equation (4.9), followed by Young's inequality on
r\^{}\{4/3\}eta\^{}\{2/3\}, gives

\[
 r=O(\eta^{1/8});\qquad
 \text{under individual calibration (4.10), }r=O(\eta^{1/4}).     \tag{4.11}
\]

These are physical annular-scale restrictions. They allow a much smaller
core scale and do not forbid all separation of scales. Positive Chern
charge alone has not been shown to produce nontrivial quotient annuli,
bounded overlap, or individual calibration. The earlier actual smooth
hyperbolic cores have trivial quotient group. The proof therefore does
not settle the unrestricted zero-end gap by itself.

\hypertarget{a-stronger-count-for-forbidden-positive-spinor-lifts}{%
\subsubsection{4.4. A stronger count for forbidden positive-spinor
lifts}\label{a-stronger-count-for-forbidden-positive-spinor-lifts}}

If a fixed marked flat lift has no invariant vector in S\^{}+ alone, the
positive-spinor part of the proof already gives

\[
 \boxed{c_1\le r^{-4}\int_{\iota Q_\Gamma}
                 (\rho-\rho_\beta)^2
                         +\int_{\iota Q_\Gamma}|F_A|^2.}        \tag{4.12}
\]

No opposite-spinor residual or torsion calibration is needed. For a
marked 1/n(1,-2), n odd, write the restricted determinant class as k in
Z/n, using the character zeta as generator. The canonical positive
weights are (0,-1); twisting by zeta\^{}a gives weights (a,a-1) and
k=2a-1. Thus a positive invariant is possible exactly when k=+1 or -1.
Every other determinant class satisfies (4.12).

The inherited amplitude estimate is sharper than the combined (4.2):
int(rho-rho\_beta)\^{}2\textless=C eta V. For bounded-overlap annuli
from a fixed finite family of such lifts, (4.3) and (4.12) give

\[
 \frac{\sum_i r_i^4}V\le C\Lambda
            [\eta+r_{\max}^4(\beta^{-1}+\sqrt\eta)]. \tag{4.13}
\]

Consequently comparable positive-density annuli have
r=O(eta\^{}\{1/4\}), without torsion calibration. More strongly, if each
annulus has r\_i\^{}\{-4\}int(rho-rho\_beta)\^{}2\textless=c\_1/2, then
int\_\{annulus\}\textbar F\_A\textbar\^{}2\textgreater=c\_1/2. Summing
and applying (4.3) to the union gives

\[
 \boxed{\frac NV\le C\Lambda(\beta^{-1}+\sqrt\eta)}  \tag{4.14}
\]

after absorbing C N r\_max\^{}4 for sufficiently small fixed r\_max.
This excludes a positive density of these calibrated forbidden lifts at
\emph{any} smaller radii; escape to arbitrarily small scales does not
evade (4.14). Its topological premise is the actual restricted Spin-c
class on the marked annulus, not just the abstract lens-space boundary.

\hypertarget{the-resonant-alternatives-solve-the-actual-local-equation}{%
\subsubsection{4.5. The resonant alternatives solve the actual local
equation}\label{the-resonant-alternatives-solve-the-actual-local-equation}}

The exceptional types in (4.6) are not merely a failure of a
representation test. They admit actual local solutions for which all
three right-hand terms in (4.8) tend to zero. This construction is on
punctured quotient balls; it makes no smooth-filling or closed-manifold
assertion.

Take hyperbolic four-space of sectional curvature -1/3 and a cyclic
group of odd order acting about its center by
(z\_1,z\_2)-\textgreater(lambda\^{}\{-1\}z\_1,lambda\^{}2z\_2), with the
lift (4.4). In rescaled normal coordinates on a fixed covering ball,
h\_r=r\^{}\{-2\}(exp\_p(r,\textbackslash cdot))\^{}*g=g\_Euc+O(r\^{}2)
in every fixed smooth norm. The radial spin identification is
equivariant. Start with

\[
 \Phi_r^0=a e_1+r c\overline z_1e_2,\qquad
 a^2=4/3,\quad c=\sqrt2/3.                          \tag{4.15}
\]

This is equivariant because bar z\_1 has character lambda. For the fixed
Clifford matrices C\_1=I, C\_2=diag(i,-i), the physical flat derivatives
are P\_1=ce\_2 and P\_2=-ice\_2, hence C\_1P\_1+C\_2P\_2=0. Direct
substitution in the cubic formula gives exactly

\[
 t_{\Phi_r^0}=ac(C_1\Phi_r^0-iC_2\Phi_r^0)
                         =2a^2c e_1,
 \qquad |t_{\Phi_r^0}|^2=128/81.                    \tag{4.16}
\]

To solve the actual equations, use the right inverses from Section 15 of
the conformal continuation. On the Euclidean covering ball set
T=(d\^{}\emph{,sqrt2 d\^{}+), R\_A=T\^{}}G\_D, R\_D=D\^{}-G\_D, where
G\_D is the Dirichlet Laplacian inverse. They are right inverses gaining
one derivative. They commute with the finite group. The O(r\^{}2)
perturbation to h\_r preserves bounded right inverses, by a Neumann
series, in fixed C\^{}\{k,alpha\} spaces. In particular no independent
incompatible boundary values for Phi or the connection are prescribed.

For a spinor in a fixed neighborhood of Phi\_r\^{}0 put

\[
 \alpha(\Phi)=R_A(0,\sqrt2 r^2\,2i\sigma_{h_r}(\Phi)),
\] \[
 \mathcal T(\Phi)=\Phi_r^0-R_D\left[
 D_{h_r,A_0}\Phi_r^0+\tfrac12c_{h_r}(\alpha(\Phi))\Phi\right].
\]

The seed Dirac error is O(r\^{}2). On a ball of radius C r\^{}2 about
that seed, this equivariant map has Lipschitz constant O(r\^{}2) and
maps the ball to itself. Its smooth fixed point satisfies
\textbar\textbar Phi-Phi\_r\^{}0\textbar\textbar{}\emph{\{C\^{}k\}+\textbar\textbar alpha\textbar\textbar{}}\{C\^{}k\}\textless=C\_k
r\^{}2. At physical scale it therefore solves

\[
 D_A\Phi=0,\qquad F_A^+=2i\sigma_g(\Phi),\qquad
 \operatorname{Ric}_g=-g,\quad W_g^\pm=0.            \tag{4.17}
\]

The actual majorant is f\_0=2, since s=-4 and W\^{}+=0. Returning to
physical derivatives gives, uniformly on the covering ball,

\[
 \rho=4/3+O(r^2),\quad t=2a^2c e_1+O(r),\quad
 \zeta=O(r),\quad |F_A|=O(1),\quad\mathcal A=O(r^2).
\]

The last A is its zero-end definition beta=infinity. For finite
beta-\textgreater infinity use a\^{}2=rho\_beta and
c\^{}2=m\_beta/(4a\^{}4), obtaining the same estimates. The actual
torsion Dirac identity, or just rescaled C\^{}1 convergence of t, shows
that the third term of (4.8) tends to zero. Thus all terms on that right
side vanish on any fixed-ratio annulus.

There is also uniform control across long local necks. Solve on an outer
cover ball of radius 4R with R-\textgreater0. The errors above are O(R)
or O(R\^{}2) throughout that ball, not errors at a subsequently chosen
inner scale. The actual torsion equation with W\^{}+=0 gives
\textbar D\_\{3A\}t\textbar\textless=C\textbar zeta\textbar\textbar P\_spin\textbar{}
=O(R). On every annulus of scale ell\textless=R, the three normalized
losses in (4.8) are respectively bounded by

\[
 C R^2,\qquad C\ell^4,\qquad C(\ell R)^{4/3}.        \tag{4.18}
\]

They tend to zero uniformly, including ell/R-\textgreater0. These are
actual local resonant necks with calibrated nonzero torsion. Their
center remains a quotient singularity when n\textgreater1; no actual
global class, closed solution, smooth filling, or boundary matching to
an external geometry has been constructed. They establish the necessary
exception in the local matching theorem, not a counterexample to Section
3 or to Yau26.

Both allowed determinant lifts have such examples. For the other lift
twist (4.4) by lambda\^{}\{-1\}. The positive representation becomes
diag(lambda\^{}\{-1\},1), and E\_- becomes diag(lambda\^{}\{-3\},1). The
seed a e\_2+r c z\_1 e\_1 is equivariant and Dirac-harmonic; its
physical cubic torsion is -2a\^{}2c e\_2. The identical contraction and
estimates apply. Thus neither of the two permitted determinant classes
can be discarded by a stronger local version of (4.12).

\hypertarget{what-an-escaping-bounded-ricci-sequence-must-carry}{%
\subsection{5. What an escaping bounded-Ricci sequence must
carry}\label{what-an-escaping-bounded-ricci-sequence-must-carry}}

The new results give a sharper description of the remaining endpoint
problem. This section is a consequence of I45 and the inherited exact
characteristic identities, not a further independent gap theorem.

Suppose an actual sequence has beta-\textgreater infinity,
eta-\textgreater0, a\^{}2\textgreater=T\_+ and uniform pointwise Ricci
bound. Write

\[
 C_+=c_2(S^+)[M]=(a^2-T_+)/4.
\]

Determinant saturation gives C\_+/V=O(beta\^{}\{-1\}+sqrt
eta)-\textgreater0, while 4pi\^{}2T\_+/V=2/3+16/(3beta\^{}2)+O(sqrt
eta). The exact rank-two formulas are

\[
 c_2(S^-)[M]=\chi+C_+,\qquad
 c_2(S^-\otimes L)[M]=\chi+2T_++9C_+.               \tag{5.1}
\]

Indeed c\_2(S\^{}-)=(a\^{}2+2chi-3tau)/4; tensoring by L adds
c\_1(S\^{}-)c\_1(L)+c\_1(L)\^{}2=2a\^{}2. It follows that

\[
 \frac{c_2(S^-\otimes L)[M]-\chi}{V}
   =\frac1{3\pi^2}+\frac8{3\pi^2\beta^2}
                     +9\frac{C_+}{V}+O(\sqrt\eta). \tag{5.2}
\]

By (3.11), chi/V must tend to infinity. Consequently both negative-
chiral bundles have divergent c\_2 density, whereas the positive-spinor
bundle has vanishing c\_2 density. The determinant twist changes this by
only a finite limiting density 1/(3pi\^{}2); it does not cause the
divergence.

Furthermore chi/T\_+-\textgreater infinity and tau/chi-\textgreater-2/3.
Since T\_+ is eventually a positive integer, chi itself tends to
infinity. In particular such a bounded-Ricci endpoint sequence cannot
remain on one fixed smooth topology.

Chern--Gauss--Bonnet also makes the curvature location explicit:

\[
 \frac1V\int|W^-|^2-\frac{8\pi^2\chi}{V}
                  =-\frac23-\frac8{3\beta^2}+O(\sqrt\eta).      \tag{5.3}
\]

Here the bounded actual model gives
avg\textbar W\^{}+\textbar\^{}2=8/(3beta\^{}2)+O(sqrt eta), avg
s\^{}2=16+O(sqrt eta), and avg\textbar Ric\^{}circ\textbar\^{}2=O(eta).
Thus the remaining bounded-Ricci escape requires unbounded opposite-Weyl
energy density, even while the full determinant curvature density stays
bounded. This explains exactly which curvature block is not paid for by
the existing positive-spinor second-order budget.

None of (5.1)--(5.3) produces individual bubbles, marked quotient
annuli, or their number. If such annuli are identified separately,
Section 4 constrains their actual lifts and costs. In particular the
local resonant examples do not by themselves account for these global
characteristic densities. The outstanding global matching problem is not
replaced by a formal count of singularities.

END COMPLETE SOURCE

\end{document}
